%% BinML methods paper -- AASTeX 7.0.1
%% Build with ./build.sh (regenerates macros + figures, then pdflatex/bibtex).
\documentclass[twocolumn,linenumbers,floatfix]{aastex701}

\input{paper_macros}   % \bml* numbers, auto-generated from canonical_numbers.json
% AUTO-GENERATED by make_gulls_macros.py -- do not edit.
\newcommand{\bmlThreshold}{0.904}
\newcommand{\bmlGapSensNone}{0.93}
\newcommand{\bmlGapSensTwo}{0.86}
\newcommand{\bmlGapSensN}{100}
\newcommand{\bmlGullsNcat}{377{,}999}
\newcommand{\bmlGullsNcatSingle}{299{,}466}
\newcommand{\bmlGullsNcatPlanet}{35{,}647}
\newcommand{\bmlGullsNcatPlanetBin}{42{,}886}
\newcommand{\bmlGullsAmpMedSingle}{0.063}
\newcommand{\bmlGullsAmpFracSingle}{41}
\newcommand{\bmlGullsSubdaySingle}{33}
\newcommand{\bmlGullsSubdayPlanet}{0.31}
\newcommand{\bmlGullsSubdayPlanetBin}{0.15}
\newcommand{\bmlGullsEligibleN}{100{,}935}
\newcommand{\bmlGullsScoredPct}{56.4}
\newcommand{\bmlGullsBetweenPct}{22.0}
\newcommand{\bmlGullsOutsidePct}{3.6}
\newcommand{\bmlGullsLowcadPct}{17.9}
\newcommand{\bmlGullsNoBaseN}{84}
\newcommand{\bmlGullsN}{56{,}975}
\newcommand{\bmlGullsNSingle}{33{,}353}
\newcommand{\bmlGullsNPlanet}{11{,}388}
\newcommand{\bmlGullsNPlanetBin}{12{,}234}
\newcommand{\bmlGullsPlanetFracW}{22}
\newcommand{\bmlGullsQmedPlanet}{1.4\times10^{-4}}
\newcommand{\bmlGullsQmedPlanetBin}{1.2\times10^{-4}}
\newcommand{\bmlGullsQksExp}{-18}
\newcommand{\bmlGullsBinSrcPct}{55}
\newcommand{\bmlGullsDetCut}{60}
\newcommand{\bmlGullsWidePass}{10}
\newcommand{\bmlGullsWideFrac}{63}
\newcommand{\bmlGullsResonantPass}{85}
\newcommand{\bmlGullsBrightSingle}{10}
\newcommand{\bmlGullsFsLowSingle}{21}
\newcommand{\bmlGullsBrightPlanet}{23}
\newcommand{\bmlGullsFsLowPlanet}{30}
\newcommand{\bmlGullsBrightPlanetBin}{21}
\newcommand{\bmlGullsFsLowPlanetBin}{35}
\newcommand{\bmlGullsRhoPnn}{0.29}
\newcommand{\bmlGullsNrhoGtOne}{43}
\newcommand{\bmlGullsFaBright}{11.8}
\newcommand{\bmlGullsFaInside}{4.1}
\newcommand{\bmlOursNonpsplQmed}{0.12}
\newcommand{\bmlOursNonpsplStellarPct}{82}
\newcommand{\bmlGullsSeasonDays}{70.7}
\newcommand{\bmlGullsPauseHoursLo}{43}
\newcommand{\bmlGullsPauseHoursHi}{44}
\newcommand{\bmlGullsGapLo}{109}
\newcommand{\bmlGullsGapHi}{120}
\newcommand{\bmlGullsObservedDays}{685}
\newcommand{\bmlGullsSpanDays}{1{,}717}
\newcommand{\bmlGullsPartialBinsPct}{34}
\newcommand{\bmlGullsBudgetLo}{2.0}
\newcommand{\bmlGullsBudgetMid}{5.2}
\newcommand{\bmlGullsBudgetHi}{11.7}
\newcommand{\bmlGullsFaShipped}{35.6}
\newcommand{\bmlGullsRecShipped}{0.214}
\newcommand{\bmlGullsRecBinShipped}{0.248}
\newcommand{\bmlGullsRecAtRecallShipped}{0.504}
\newcommand{\bmlGullsRhoLoShipped}{0.402}
\newcommand{\bmlGullsRhoHiShipped}{0.258}
\newcommand{\bmlGullsFaGapaware}{11.6}
\newcommand{\bmlGullsRecGapaware}{0.295}
\newcommand{\bmlGullsRecBinGapaware}{0.313}
\newcommand{\bmlGullsRecAtRecallGapaware}{0.457}
\newcommand{\bmlGullsRhoLoGapaware}{0.046}
\newcommand{\bmlGullsRhoHiGapaware}{0.674}
\newcommand{\bmlGullsFaFspl}{4.8}
\newcommand{\bmlGullsRecFspl}{0.390}
\newcommand{\bmlGullsRecBinFspl}{0.421}
\newcommand{\bmlGullsRecAtRecallFspl}{0.378}
\newcommand{\bmlGullsRhoLoFspl}{0.031}
\newcommand{\bmlGullsRhoHiFspl}{0.288}
\newcommand{\bmlGullsFaFsplOne}{5.2}
\newcommand{\bmlGullsRecFsplOne}{0.369}
\newcommand{\bmlGullsRecBinFsplOne}{0.399}
\newcommand{\bmlGullsRecAtRecallFsplOne}{0.370}
\newcommand{\bmlGullsRhoLoFsplOne}{0.032}
\newcommand{\bmlGullsRhoHiFsplOne}{0.333}
\newcommand{\bmlGullsFaFsplFive}{8.7}
\newcommand{\bmlGullsRecFsplFive}{0.344}
\newcommand{\bmlGullsRecBinFsplFive}{0.372}
\newcommand{\bmlGullsRecAtRecallFsplFive}{0.445}
\newcommand{\bmlGullsRhoLoFsplFive}{0.040}
\newcommand{\bmlGullsRhoHiFsplFive}{0.464}
\newcommand{\bmlGullsFaNoisy}{12.1}
\newcommand{\bmlGullsRecNoisy}{0.298}
\newcommand{\bmlGullsRecBinNoisy}{0.320}
\newcommand{\bmlGullsRecAtRecallNoisy}{0.475}
\newcommand{\bmlGullsRhoLoNoisy}{0.053}
\newcommand{\bmlGullsRhoHiNoisy}{0.513}
\newcommand{\bmlGullsFaOnset}{5.3}
\newcommand{\bmlGullsRecOnset}{0.375}
\newcommand{\bmlGullsRecBinOnset}{0.407}
\newcommand{\bmlGullsRecAtRecallOnset}{0.381}
\newcommand{\bmlGullsRhoLoOnset}{0.031}
\newcommand{\bmlGullsRhoHiOnset}{0.326}
\newcommand{\bmlGullsNcheckpoints}{6}
\newcommand{\bmlGullsBudgetSlack}{0.05}
\newcommand{\bmlGullsColourGainLo}{0.042}
\newcommand{\bmlGullsColourGainHi}{0.068}
\newcommand{\bmlGullsTruthN}{6{,}600}
\newcommand{\bmlGullsTruthNSingle}{1{,}600}
\newcommand{\bmlGullsTruthNPlanet}{2{,}500}
\newcommand{\bmlGullsEpochExcess}{11}
\newcommand{\bmlGullsUndetPlanet}{44}
\newcommand{\bmlGullsUndetPlanetBin}{46}
\newcommand{\bmlGullsDetPlanet}{56}
\newcommand{\bmlGullsDetPlanetBin}{54}
\newcommand{\bmlGullsUndetPlanetW}{41}
\newcommand{\bmlGullsFailChi}{26}
\newcommand{\bmlGullsFloorVeto}{19}
\newcommand{\bmlGullsDetRescaled}{55}
\newcommand{\bmlGullsSinglePsplRefitPct}{2.1}
\newcommand{\bmlGullsRecDetPlanet}{0.474}
\newcommand{\bmlGullsRecDetPlanetBin}{0.557}
\newcommand{\bmlGullsRecGenPlanet}{0.380}
\newcommand{\bmlGullsRecGenPlanetBin}{0.404}
\newcommand{\bmlGullsFlagUndet}{24}
\newcommand{\bmlGullsFlagFloorVeto}{40}
\newcommand{\bmlGullsSingleMisfitK}{10}
\newcommand{\bmlGullsSingleMisfitN}{20}
\newcommand{\bmlGullsSingleMisfitLo}{30}
\newcommand{\bmlGullsSingleMisfitHi}{70}
\newcommand{\bmlGullsFloorFlagThreeLo}{38}
\newcommand{\bmlGullsFloorFlagThreeHi}{47}
\newcommand{\bmlGullsFloorFlagOneLo}{16}
\newcommand{\bmlGullsFloorFlagOneHi}{20}
\newcommand{\bmlGullsRecChiLo}{0.17}
\newcommand{\bmlGullsRecChiHi}{0.87}
\newcommand{\bmlGullsColDetOne}{0.534}
\newcommand{\bmlGullsColDetThree}{0.539}
\newcommand{\bmlGullsColUndetOne}{0.170}
\newcommand{\bmlGullsColUndetThree}{0.270}
\newcommand{\bmlGullsPrecSample}{0.694}
\newcommand{\bmlGullsPrecScored}{0.618}
\newcommand{\bmlSidecarRecDetPlanet}{0.325}
\newcommand{\bmlSidecarRecDetPlanetBin}{0.398}
\newcommand{\bmlGullsFloorDetLo}{70}
\newcommand{\bmlGullsFloorDetHi}{40}
\newcommand{\bmlGullsSubdayN}{1{,}385}
\newcommand{\bmlGullsSubdayTe}{0.13}
\newcommand{\bmlGullsSubdayFa}{0.7}
\newcommand{\bmlGullsSubdayFaW}{3.1}
\newcommand{\bmlGullsSubdayPerLo}{60}
\newcommand{\bmlGullsSubdayPerHi}{93}
\newcommand{\bmlGullsNoiseBright}{0.95}
\newcommand{\bmlGullsNoiseFaint}{2.64}
\newcommand{\bmlGullsNoiseBkg}{7}
\newcommand{\bmlGullsNoiseGlobal}{1.41}
\newcommand{\bmlSidecarThr}{0.957}
\newcommand{\bmlSidecarThrLo}{0.942}
\newcommand{\bmlSidecarThrHi}{0.962}
\newcommand{\bmlSidecarFa}{1.6}
\newcommand{\bmlSidecarRecPlanet}{0.24}
\newcommand{\bmlSidecarRecPlanetBin}{0.27}
\newcommand{\bmlSidecarFaLo}{1.3}
\newcommand{\bmlSidecarFaHi}{2.4}
\newcommand{\bmlSidecarPoolPrev}{6.0}
\newcommand{\bmlSidecarPoolCompl}{0.677}
\newcommand{\bmlOccN}{6{,}024}
\newcommand{\bmlOccFaObs}{4.9}
\newcommand{\bmlOccFaOne}{3.6}
\newcommand{\bmlOccRecObs}{0.35}
\newcommand{\bmlOccRecOne}{0.26}
    % \bmlGulls*/\bmlGap*/\bmlCg*/... numbers for the RMDC26 section, from make_gulls_macros.py (fail-closed)

\newcommand{\code}[1]{\texttt{#1}}
\newcommand{\dchisq}{\ensuremath{\Delta\chi^2}}
\newcommand{\Rom}{\textit{Roman}}

\shorttitle{BinML: partial-light-curve anomaly triage}
\shortauthors{Bhatia}

\begin{document}

\title{BinML: Detectability-Conditioned Anomaly Triage for Synthetic \\
\textit{Roman}-like Partial Light Curves}

\author[0009-0007-4447-6325]{Kunal Bhatia}
\affiliation{Independent Researcher}
\email{kunal29bhatia@gmail.com}

\begin{abstract}
\Rom{}'s Galactic Bulge Time-Domain Survey will yield more microlensing light curves than anyone can inspect,
whose anomalies must be found while the events are in progress. We present \textbf{BinML}, a
$\bmlNparams$-parameter network that classifies partial \Rom{}-like light curves in three bands into six classes,
among them detectable binary-lens anomalies. Two choices make the task well posed: labels are conditioned on an
adopted detectability policy, so a binary whose perturbation falls below its thresholds is labelled a single lens,
and on the revealed part of the season, so the label follows the cascade flat $\rightarrow$ single lens
$\rightarrow$ anomaly as evidence arrives. On complete synthetic seasons BinML recovers $\bmlCompletenessPct\%$ of
detectable anomalies at purity $\bmlPurity$. Revealed in half-day steps in F146, $\bmlCascDetFrac\%$ of eligible
binaries alert within the season, $\bmlCascPrematurePct\%$ before a truth-informed onset and the rest a median
$\bmlCascMedianLagNonPrem$~days after it; a three-band stream of every class yields
$\bmlRefStreamAlertsDay$ alerts per 1,000 events per day, $\bmlRefStreamPurity\%$ of them detectable anomalies. On
RMDC26, an independent GULLS simulation of the survey, the shipped model fails, flagging $\bmlGullsFaShipped\%$ of
single lenses, mostly because our training seasons lack the survey's F146 observing pauses. Fine-tuning on the
measured pauses with finite-source single lenses brings that to $\bmlGullsFaCombined\%$ and raises planetary recall
at a $\bmlGullsBudgetMid\%$ false-alarm budget from $\bmlGullsRecShipped$ to $\bmlSeedRecLo$--$\bmlSeedRecHi$ (three
seeds); RMDC26 guided the checkpoint choice, so these are optimistic. All are synthetic benchmarks
($0.02$~mag detectability floor, $\bmlPriorPiZero\%$ anomaly prevalence), not yield forecasts. Code,
weights and artifacts are public.
\end{abstract}

\keywords{Gravitational microlensing (672) --- Exoplanet detection methods (489) ---
Neural networks (1933) --- Time domain astronomy (2109) --- Convolutional neural networks (1938)}

\section{Introduction}\label{sec:intro}

Gravitational microlensing complements other planet-detection methods because it responds to a
lens's mass rather than its light. It is sensitive to cold, low-mass planets far from their host
stars and to unbound planetary-mass objects drifting through the Galaxy---populations that transit
and radial-velocity surveys, biased toward hot and close-in planets, recover less readily
\citep{Mao1991,GouldLoeb1992,Suzuki2016}. The Nancy Grace \Rom{} Space Telescope will exploit
this sensitivity systematically. Its Galactic Bulge Time-Domain Survey (GBTDS) will monitor the
crowded bulge fields at high cadence in the near-infrared. It is predicted to detect of order a
thousand bound planets and to constrain the free-floating planet mass function
\citep{Spergel2015,Penny2019,Johnson2020,Sumi2023,Gaudi2012}, and the same data will yield
transiting planets \citep{Wilson2023}. In its design
\citep{ROTAC2025,STScI2025GBTDS}, the wide F146 filter samples the light curve every
$\sim12$~minutes; two auxiliary bands, F087 and F213, add colour information on staggered
$\sim$6-hour cycles.

That sensitivity brings an operational problem of scale. Genuine microlensing events are a small
minority of the survey light curves, detectable binary-lens anomalies are rarer still, and fitting
a full binary-lens model to every source is impractical. A cheap classifier can therefore rank
which partial light curves deserve inspection or detailed modelling. Timing matters to that task:
before a caustic perturbation appears, a binary lens can resemble an ordinary single lens. A model
evaluated only on complete seasons can assign an anomalous label before the evidence that defines
it has entered the observed window. Existing rapid-response systems already show the value of
low-latency anomaly detection in microlensing \citep{Udalski2003,Dominik2007,Dominik2008,Tsapras2009}.
Here we ask a narrower question: as an eligible synthetic binary is progressively revealed, when
does the classifier first cross an anomaly threshold relative to a truth-informed onset?

Machine learning has already been used at the scale relevant to microlensing surveys.
\citet{Wyrzykowski2015} searched 150 million OGLE-III sources and classified 3,718 candidates, while
\citet{Husseiniova2021} applied a machine-learning search to roughly 700 million VVV light curves;
both were completed-data searches. The random-forest microlensing classifier of \citet{Godines2019} and neural networks on
ground-based photometry \citep{Mroz2020} address related classification tasks. KMTNet's completed-event finder \citep{Kim2018EventFinder} and the subsequent
\textsc{AnomalyFinder} analyses \citep{Zang2022}, which search residuals from fitted single-lens
models, demonstrate systematic event and anomaly discovery in real survey data. For simulated high-cadence \Rom{} light curves,
\citet{Khakpash2021} developed morphology and model-fit features intended for rapid classification.
In the alert-broker setting, Fink's LIA module
performs real-time microlensing classification but does not include a binary-anomaly class
\citep{Moller2021}. For \Rom{}, \citet{Zhang2021} demonstrated real-time posterior inference for a
known binary-lens event; that is downstream parameter inference, not discovery in the full stellar
stream. Most closely, \citet{CrispimRomao2025} evaluated an isolation-forest alarm on drip-fed
synthetic LSST light curves containing microlensing, binary lenses, and variables. Their online
stage separates signal from constant sources and their physical classification is offline; our
contribution is a supervised six-class \Rom{}-like model with time-dependent labels. Related
time-domain work also studies early classification of explosive transients
\citep{Muthukrishna2019}. Architecturally we use standard time-series transformer components
\citep{Astromer2023,Vaswani2017}, and our evaluation shares the prior-shift concern highlighted by
PLAsTiCC \citep{Hlozek2020}.

Across this work we find no trained classifier that combines all of the elements BinML needs:
multi-band \Rom{}-cadence input, a binary and planetary anomaly class evaluated alongside the
variable-star populations that contaminate it, operation on partial rather than completed light
curves, and released weights. The searches above are single-band, completed-data, or
binary-blind; the \Rom{}-specific work stops at feature definition or at downstream inference
for an already-identified event. To our knowledge BinML is the first released model to occupy
that combination, which is what the rest of this paper evaluates.

This paper makes four contributions. First, we introduce \emph{detectability-conditioned
labelling} (\S\ref{sec:sims}): events are labelled by whether they cross adopted synthetic
thresholds, aligning the training target with that synthetic policy.
Second, we build a \emph{partial-light-curve cascade} (\S\ref{sec:cascade}) by making each label a
function of how much of the season has been revealed. Third, we evaluate the model at synthetic
survey scale (\S\ref{sec:eval}--\ref{sec:limits}):
completeness at a fixed purity as the headline metric, a selection-aware reweighting that partially corrects
for how the training set was sampled, and stress tests on an independent same-prior sample, targeted regimes and out-of-range sweeps (a
$\bmlStressN$-million-event collection scored with the released model's predecessor, and a regenerated subset
scored with both). Fourth, we test transfer to RMDC26,
an independent simulation of the survey made with a different simulator (\S\ref{sec:gulls}): it
exposes a failure of the shipped model on the survey's observing schedule, which a fine-tune with gap
augmentation reduces.

\section{Simulated survey and labelling}\label{sec:sims}

\begin{figure*}[tbp]
\centering
\includegraphics[width=0.98\textwidth]{outputs/figures/lightcurve_gallery.pdf}
\caption{Illustrative simulated \Rom{}-like light curves, one per class (F146 in grey, the
bluer F087 in blue, the redder F213 in red), under the simulator assumptions. Only bands with
surviving epochs are plotted (the periodic-variable example has no F087 epochs), and heavy dust leaves F087
undetected for $\bmlNoFilterBlue\%$ of detectable-anomaly events. \texttt{Flat} sources are
noise-dominated; \texttt{PSPL} shows the smooth single-lens magnification; \texttt{NonPSPL}
carries the sharp caustic features that mark a binary or planetary lens. The simulated lens
magnification is intrinsically achromatic, although blending makes the observed amplitudes
band-dependent. The three
variable-star classes span short-period pulsators and eclipsing binaries (\texttt{PeriodicVar},
shown phase-folded), slow red-giant variability (\texttt{LongPeriodVar}), and outbursts
(\texttt{Eruptive}). The visual separability of the caustic signature from ordinary single-lens
curvature is the information the model must preserve. The examples are deliberately selected for
visual legibility and correct model classification rather than sampled at random; they are not
performance evidence. Aggregate results are reported in \S\ref{sec:results}.}
\label{fig:gallery}
\end{figure*}

\subsection{Six classes across three bands}
We simulate \Rom{} GBTDS light curves in the six selected classes of Figure~\ref{fig:gallery}.
\texttt{Flat} sources show no
detectable event. \texttt{PSPL} events are single-lens, point-source point-lens microlensing
\citep{Paczynski1986}. \texttt{NonPSPL} events---the anomalous, science-critical class---are
binary and planetary lenses, computed with the contour-integration code \textsc{VBBinaryLensing}
\citep{Bozza2010,Bozza2018}, and are distinguished by the caustic perturbations visible in the
third panel of Figure~\ref{fig:gallery}. The remaining three classes are the variable-star
populations that dominate the real bulge and are the classifier's most dangerous contaminants:
\texttt{PeriodicVar} (RR~Lyrae, eclipsing binaries, $\delta$~Scuti), \texttt{LongPeriodVar}
(Miras, semiregulars, and OSARGs---small-amplitude red-giant pulsators---whose slow rise most
readily mimics a long-timescale microlensing event), and \texttt{Eruptive} (dwarf novae and Be
outbursts). We adopt an authored truncated log-normal event-timescale distribution whose mean is
anchored to an empirical inner-bulge estimate \citep{Mroz2019}. Binary mass ratio and projected
separation are instead drawn log-uniformly over broad analytic supports; they are not a fitted
population model \citep{Suzuki2016,Zhang2021}. Variable-star periods and amplitudes are motivated
by OGLE bulge catalogues, but the generated curves are phenomenological analytic waveform families
rather than sampled empirical templates \citep{Soszynski2013,Soszynski2014}. Every event is realised
on a hybrid \Rom{}-like schedule: legacy-like $\sim$15-minute wide-band sampling and a single
72-day season, combined with 6-hour visits in each of the current F087 and F213 auxiliary filters.
The curves are degraded by photometric noise, blending, and interstellar extinction. The released
model was trained on this legacy schedule (\S\ref{sec:limits}).

The source population is likewise parametric rather than a Galactic population synthesis. The
observed F146 baseline is uniform from 20 to 25~mag and $A_{K_s}$ is uniform from 0.2 to 1.2,
sampled independently. Baseline colour offsets come from a grey intrinsic source spectrum plus the
adopted extinction law; intrinsic stellar colour and its correlations with luminosity, variability,
and line of sight are not modelled. These choices give a classifier benchmark broad coverage.

Table~\ref{tab:simulation} collects the assumptions most likely to affect transfer.

\begin{table*}[tbp]
\centering
\caption{Core simulation and labelling assumptions for the released model.\label{tab:simulation}}
\small
\hrule
\vspace{4pt}
\begin{minipage}[t]{0.47\textwidth}
\textbf{$t_E$.} Truncated log-normal, median 16.3 d,
$\sigma_{\log_{10}}=0.35$, 1--300 d; authored width, with the mean anchored to
\citet{Mroz2019}.\par\smallskip
\textbf{$u_0$.} $\mathcal{U}(0,2)$; broad analytic support following
\citet{Zhang2021}.\par\smallskip
\textbf{$q$; $s$.} Log-$\mathcal{U}(10^{-6},1)$; log-$\mathcal{U}(0.2,5)$;
benchmark supports, not occurrence rates.\par\smallskip
\textbf{$\rho$; $\alpha$.} Log-$\mathcal{U}(10^{-4},10^{-2})$;
$\mathcal{U}(0,2\pi)$; binary-lens geometry support.\par\smallskip
\textbf{$f_s$.} Log-$\mathcal{U}(0.1,1)$; reference-band source-flux
fraction.\par\smallskip
\textbf{$t_0$.} $\mathcal{U}[-p,72+p]$ d, where
$p=\min(0.5t_E,36\,{\rm d})$; permits partial and out-of-season peaks.
\end{minipage}\hfill
\begin{minipage}[t]{0.47\textwidth}
\textbf{F146 baseline.} $\mathcal{U}(20,25)$ AB mag; observed reference-band
baseline.\par\smallskip
\textbf{$A_{K_s}$; intrinsic SED.} $\mathcal{U}(0.2,1.2)$ mag; grey in AB;
colours arise only from extinction and blending.\par\smallskip
\textbf{Cadence and window.} F146 every 15 min; F087 and F213 every 6 h; one
72 d season; legacy design.\par\smallskip
\textbf{Noise.} Flux-space source-plus-background noise and a 0.001 mag floor,
sampled independently at each epoch.\par\smallskip
\textbf{Label policy.} $\Delta\chi^2_{\rm event}\geq500$,
$\Delta\chi^2_{\rm anom}\geq160$, and amplitude $\geq0.02$ mag; noise-free
statistics and an unvalidated amplitude floor.
\end{minipage}
\vspace{4pt}
\hrule
\vspace{3pt}

\parbox{0.94\textwidth}{\footnotesize \textit{Note.} Log-$\mathcal{U}$ denotes a distribution
uniform in the logarithm. Variable-star mixtures use the phenomenological waveform families
described in the text and released generator; class quotas are deliberately selected for training
and are not a survey event-rate model.}
\end{table*}

The \texttt{NonPSPL} generator contains static binary lenses (2L1S) with finite-source
magnification, but not binary-source (1L2S) events, annual parallax, or lens orbital motion. The
1L2S omission is the most consequential of the three: a binary source produces a bump on an
otherwise single-lens curve that can mimic a binary-lens anomaly, and it is a long-recognised
degeneracy with planetary-mass 2L1S perturbations \citep{Gaudi1998}. It is therefore the dominant
astrophysical false positive that our purity estimate does not contain, and adding it would
lower the reported \texttt{NonPSPL} purity by an amount we have not measured. Annual parallax
and lens orbital motion distort the light-curve shape but do not manufacture a caustic-like
feature, so we expect them to matter less over a single 72-day season. The contaminant set
likewise omits many real irregular and correlated systematics.

\subsection{Labelling by an adopted synthetic detectability policy}\label{sec:detcond}
BinML's policy-conditioned synthetic labels are set by explicit thresholds applied to simulated,
noise-free curves, rather than by generator class alone; we use ``detectability-conditioned'' as
shorthand for this policy. A microlensing event whose in-window amplitude falls below the adopted
floor is labelled \texttt{Flat}. Likewise, a binary whose caustic perturbation falls below the
anomaly thresholds is labelled \texttt{PSPL}. For binaries we measure the departure of the
noiseless curve from a refitted single-lens model and withhold the anomalous label when that
departure is too small under the policy.

The class definitions---and therefore every metric in this paper---depend on these thresholds. An
event must reach $\Delta\chi^{2}_{\mathrm{event}} \ge 500$ against a flat baseline to be a
detected event at all; a binary must reach $\Delta\chi^{2}_{\mathrm{anom}} \ge 160$ against its
best-fitting single-lens model to keep the anomalous label; and both statistics must exceed a peak
amplitude of $0.02$~mag. Both $\Delta\chi^{2}$ values are computed on noise-free curves, so neither
is a likelihood-ratio statistic: under the null each is identically zero, and each grows in
proportion to the epoch count. At \Rom{}'s sampling that makes both $\chi^{2}$ cuts close to
vestigial for our error model and leaves the $0.02$~mag amplitude floor as the operative criterion. That floor is an
adopted simulation policy rather than a measured instrumental limit. Section~\ref{sec:ablations}
measures what the labelling scheme buys against the alternative of labelling by generator intent.

The scale of the relabelling is the point. Across the stored events of the final test set,
$\bmlRelabelled\%$ receive a policy-conditioned label that differs from the class they were generated as. The effect is
concentrated where it should be. Of the stored events generated as binary or planetary lenses, only
$\bmlNonpsplKept\%$ show a detectable anomaly and keep the anomalous label; $\bmlNonpsplToPspl\%$
have a caustic perturbation below those thresholds and are labelled as single lenses. (Those
single-lens by-products are stored only in part, \S\ref{sec:eval}; over all generated binaries the two
fractions are $\bmlNonpsplKeptW\%$ and $\bmlNonpsplToPsplW\%$.) Similarly,
$\bmlLpvToFlat\%$ of the long-period variables vary too weakly within a single season to be
detected and are labelled flat. Training against generator
intent asks the network to assign an anomalous or variable label on about a third of its examples
even though the corresponding signal falls below the adopted policy threshold.

\section{Model}\label{sec:model}

\begin{figure*}[tbp]
\centering
\includegraphics[width=0.92\textwidth]{outputs/figures/schematic.pdf}
\caption{BinML architecture. Each of the three \Rom{} bands is binned onto a fixed grid whose
bins carry five channels---baseline-relative mean, min, and max magnitude, observed fraction, and
an observed mask. A convolutional stem, whose non-learned min/max lanes preserve caustic extrema
by construction, reduces the bands to $\bmlNtokens$ tokens that a four-layer transformer
classifies into six classes. The whole model has $\bmlNparams$ parameters.}
\label{fig:schematic}
\end{figure*}

BinML has two parts (Figure~\ref{fig:schematic}). The first is a small convolutional network---a
stack of learned filters that scan along the light curve looking for local shapes, the same kind
of building block used in image recognition. The second is a transformer, a now-standard
architecture that lets every part of the light curve ``attend to'' every other part, so the model
can weigh, for example, a brief brightening against the overall trend around it. Together they
come to just $\bmlNparams$ trainable numbers, which by the standards of modern deep learning is
tiny. The small size is deliberate. We are not short of training data---we simulate millions of events
($\bmlTrainFinalEvents$ million in the final training stage alone)---so the pressure is not to build a larger model but to build one cheap enough to
run on the entire survey, night after night. A compact network that classifies a light curve in a
single pass is far more useful here than a heavy one applied only to a pre-selected handful of
candidates.

The model reads the light curve in a form that keeps the survey's own timing. Each of the three
filters is divided into a fixed set of time bins across the 72-day observing season: the workhorse
F146 filter into $\bmlBinsF$ bins (from $\bmlEpochsF$ individual exposures taken every
$\bmlCadence$~minutes), and the two colour filters into 96 bins each. For every bin the model is
given five simple numbers: the average brightness in that bin, the brightest and faintest points
in it, the fraction of the bin that was actually observed, and a flag saying whether it was
observed at all. All brightnesses are measured relative to a single quiet-state (baseline)
level defined in F146 and subtracted from all three filters. This removes absolute source
brightness while retaining the simulated inter-band offsets, including the effects of extinction
and blending, but not an independently measured intrinsic source colour.

The training and frozen evaluations use the simulator-provided F146 baseline. A deployed pipeline
would estimate that level from a partial, possibly magnified light curve; the released inference
interface includes an estimator, whose error we have not evaluated against the generator value.

Keeping the brightest and faintest point in each bin, rather than only the average, retains narrow
caustic structure during time binning. In the largest measured case, mean pooling
would reduce a nearly three-magnitude excursion to about a third of a magnitude. We therefore
propagate exact extrema through each strided reduction stage before a learned projection combines
the channels; the extrema are not supplied unchanged to the final classifier. When a colour filter
is missing for an event---the
bluer F087 filter yields no usable epoch for $\bmlNoFilterBlue\%$ of detectable-anomaly
events behind heavy dust---the model is simply told to
ignore it rather than being fed a guess. A final pooling step combines the processed bins into a
single summary, from which a standard classifier reads off the six class probabilities. The model
uses only the binned brightnesses; it does not use the individual error bars.

This compactness makes repeated survey-scale scoring practical. Classification is a single forward
pass with no per-event fitting. On a $\bmlInferCores$-core \bmlInferCpu{} CPU with
$\bmlInferThreads$ threads and otherwise idle, the model scores $\bmlInferEps$ pre-tokenised light curves per second
($\bmlInferMs$~ms per event) in batches of 1024; the rate falls when other work loads the machine.
\code{validation/inference\_benchmark.py} regenerates this together with the processor, thread count, load average
and timing protocol, and replaces the figure of an earlier revision, whose hardware and script were not retained.

\section{Partial-light-curve evaluation}\label{sec:cascade}

\begin{figure*}[tbp]
\centering
\includegraphics[width=0.98\textwidth]{outputs/figures/cascade_evolution.pdf}
\caption{The partial-light-curve cascade on three simulated events as the season is revealed. \emph{Top:} the
light curve in the bands with surviving epochs (F146 grey, F087 blue, F213 red); the dashed line marks the anomaly
onset $t_{\mathrm{anom}}$. \emph{Bottom:} BinML's class probability against the days of the season revealed.
The clean binary moves \texttt{Flat}$\rightarrow$\texttt{PSPL}$\rightarrow$\texttt{NonPSPL}, jumping to
\texttt{NonPSPL} near the onset rather than during the smooth rise; the subtle binary follows the same order with a
later, weaker anomaly; the single-lens control acquires no appreciable \texttt{NonPSPL} probability. The three
events are selected illustrations, not evidence; the scan below measures the first \texttt{NonPSPL} threshold crossing,
not the full three-state sequence.}
\label{fig:cascade}
\end{figure*}

A streaming classifier receives a partial light curve. Before its anomaly appears a binary lens is photometrically
a single lens, so assigning the complete-season class to every prefix would teach the network a target the revealed
window does not support. We therefore make the label a function of observation time. For each binary we compute the
day $t_{\mathrm{anom}}$ at which its anomaly first crosses the adopted thresholds, from the noise-free light curve
and a single-lens fit seeded at the injected parameters. During training the network sees light curves truncated at
random days: a binary truncated before $t_{\mathrm{anom}}$ carries the \texttt{PSPL} label, with the detectability
floor re-evaluated on the truncated data, and earns \texttt{NonPSPL} once the anomaly lies inside the revealed
window. The training labels resolve $t_{\mathrm{anom}}$ only to the end of the $7.2$-day interval in which the
anomaly became detectable, a conservative bias that usually delays the label (the evaluation onsets below are
resolved at $0.5$~days). The prefix labels of the other classes are also generator-informed: an eruptive source
changes class at its injected start time, and periodic variability uses an analytic amplitude--period proxy rather
than the realised waveform. Measured against the label rule itself on one training shard ($\bmlTruthN$ events,
$\bmlTruthReps$ truncations each, with the noise-free signal stored), the released training labels differ from it
for $\bmlTruthTruncNonLegacy\%$ of truncated binary presentations, about half of them still \texttt{PSPL} after the
anomaly had become detectable, and for $\bmlTruthTruncLpv\%$, $\bmlTruthTruncPer\%$, $\bmlTruthTruncPspl\%$ and
$\bmlTruthTruncErupt\%$ of truncated long-period, periodic, single-lens and eruptive presentations; with the stored
signal and a half-day onset the binary disagreement falls to $\bmlTruthTruncNonFixed\%$. The released model is
trained with the approximate labels.

Figure~\ref{fig:cascade} shows the cascade on individual events, and Figure~\ref{fig:probconf} extends it to a clear
and a marginal event of every class, each re-observed 30 times with fresh photometric noise: the clear events settle
once their evidence appears, whereas the marginal ones wander as the season is revealed, and for most of them the
noise draw alone moves the end-of-season call by tens of percentage points; the flat source stays near 1.

\begin{figure*}[tbp]
\centering
\includegraphics[width=0.96\textwidth,height=0.74\textheight,keepaspectratio]{outputs/figures/prob_evolution_confidence.pdf}
\caption{Class probability at every half-day cut for one clear and one marginal event per class (released model):
light curve left (insets fold the periodic variables at the true period), class probabilities right. Lines are
medians and bands the 16--84\% range over 30 re-observations of the same event with fresh photometric noise; the
numbers give the event's own-class probability at season end. The clear event is the first of its class in seed
order with a legible signal that is classified correctly on the complete season (for \texttt{PSPL}, peaking inside
the season); the marginal one is, among the first 300 of its class, the one whose end-of-season own-class
probability stays closest to 0.5 over repeated noise draws; for \texttt{Flat} even that event is called with
probability near 1. Dotted: the \texttt{NonPSPL} operating threshold; dashed: the anomaly onset (generator grid). Generated by
\code{paper/make\_data\_figures.py}, which checks these statements on the events it draws.}
\label{fig:probconf}
\end{figure*}

\paragraph{Protocol.} An \emph{eligible event} is a simulated binary whose anomaly survives
detectability-conditioned labelling and has a finite onset; eligibility never depends on the model. For each of
$\bmlCascN$ eligible binaries we reveal the season in half-day steps, scoring F146 alone, and record the
\emph{alert}, the first cut at which $P(\texttt{NonPSPL})$ reaches the operating threshold; an event with no alert
is right-censored. An alert is \emph{premature} if it precedes the onset, and the \emph{lag} is alert time minus
onset time. The onset is ambiguous, because the anomaly statistic is not monotone in revealed time: as more of the
season arrives the best single-lens fit can absorb an early deviation, which happens for $\bmlCascNonMonotone$ of
the $\bmlCascN$ events. We therefore bracket it, between the \emph{first-detectable} cut and the cut after which
the anomaly stays detectable to the end of the season (\emph{persistently detectable}), and report both; the
generator's own $7.2$-day grid rounds the onset upward by a median of $\bmlCascOnsetShift$~days and counts on-time
alerts as premature.

\paragraph{Results.} BinML alerts on $\bmlCascDetFrac\%$ of eligible binaries within the season;
$\bmlCascCensored$ are right-censored. The alert precedes the onset for $\bmlCascPrematurePct\%$ of eligible
events (95\% CI $[\bmlCascPrematureLo,\bmlCascPrematureHi]\%$) under the first-detectable onset and
$\bmlCascPrematurePersistPct\%$ under the persistent one, and the median lag of the other alerts is
$\bmlCascMedianLagNonPrem$ or $\bmlCascMedianLagPersist$~days (Figure~\ref{fig:cascade_summary}). Because the scan
is stored, the decision rule can be varied on the same events (Table~\ref{tab:policy}): a daily or two-day grid
lowers the premature rate to $\bmlCascGridOnePct$ and $\bmlCascGridTwoPct\%$ at almost unchanged detection, at the
price of later alerts; requiring two consecutive crossings gives $\bmlCascPersistTwoPct\%$ premature at
$\bmlCascPersistTwoDet\%$ detection; revealing all three bands lowers premature alerts to $\bmlCascMultibandPct\%$
and raises detection to $\bmlCascMultibandDet\%$ but lengthens the lag to $\bmlCascMultibandMedianLag$~days; and
alerting on the argmax class instead of the threshold makes detection near-complete at $\bmlCascArgmaxPct\%$
premature. The timing is strongly policy-dependent. Whether the truncation augmentation \emph{causes} the intended
ordering is tested in \S\ref{sec:ablations}.

\begin{deluxetable*}{lcccc}
\tablecaption{Premature-alert rate and within-season detection under different onset definitions and alert
policies, on the same $\bmlCascN$ eligible binaries from one stored scan. Intervals are Wilson $95\%$ on the
eligible denominator; ``median lag'' is over alerts that were not premature. Generated from
\code{validation/cascade\_reproduce\_result.json}.\label{tab:policy}}
\tablehead{\colhead{variation} & \colhead{premature (\% eligible)} & \colhead{95\% CI (\%)} &
           \colhead{alerted (\%)} & \colhead{median lag (d)}}
\startdata
\input{outputs/cascade_policy_table.tex}
\enddata
\tablenotetext{a}{The primary protocol: F146 only, half-day grid, single threshold crossing, first-detectable onset.
Every other row changes one of those choices.}
\end{deluxetable*}

\begin{figure}[tbp]
\centering
\includegraphics[width=0.9\columnwidth]{outputs/figures/cascade_summary.pdf}
\caption{Cumulative first-alert time for the $\bmlCascN$ eligible binaries, F146 only, half-day grid. The height at
$t<0$ (shaded) is the premature-alert fraction, the plateau the within-season detection fraction, the shortfall to
1 the right-censored events. Solid: alerts scored against the half-day onset; dashed: against the generator's coarse
$7.2$-day onset, which inflates the premature fraction. From \code{validation/cascade\_trace.npz}.}
\label{fig:cascade_summary}
\end{figure}

\paragraph{Every class in the stream.} The scan above says nothing about how many alerts a stream would raise, so
we also revealed every event of the two held-out-pool shards of \S\ref{sec:limits} ($\bmlRefStreamN$ events of all
six classes; $\bmlRefValPct\%$ of them are rows on which the frozen threshold was chosen) in the same half-day
steps, with all three bands and with F146 alone. With three bands no flat source or synthetic variable raised an
alert; $\bmlRefStreamAlertPsplW\%$ of single lenses (reweighted) and $\bmlRefStreamAlertNonpspl\%$ of
\texttt{NonPSPL} events did, which at the simulated mixture ($\bmlRefStreamPrev\%$ anomaly prevalence) is
$\bmlRefStreamAlertsDay$ alerts per 1,000 events per day, $\bmlRefStreamPurity\%$ of them from detectable anomalies
($\bmlRefStreamPurOne\%$ at a 1\% prevalence and $\bmlRefStreamPurTenth\%$ at 0.1\%, by the prior-shift arithmetic
of Figure~\ref{fig:prior}). Nearly all single-lens alerts ($\bmlRefStreamDemoted\%$ after reweighting) come from
binaries whose anomaly falls below the detectability policy; only $\bmlRefStreamGenK$ of the $\bmlRefStreamGenN$
generated single lenses with a \texttt{PSPL} label alerted. With F146 alone the stream raises
$\bmlRefStreamFAlertsDay$ alerts per 1,000 events per day, $\bmlRefStreamFPurity\%$ of them detectable anomalies,
and $\bmlRefStreamFVarAlerts$ flat sources or variables alert. Among the \texttt{NonPSPL} events
$\bmlRefStreamPrem\%$ alerted before the half-day onset with three bands and the median lag was
$\bmlRefStreamLag$~days, close to the three-band row of Table~\ref{tab:policy}.

\section{Evaluation methodology}\label{sec:eval}

Three choices define what the synthetic benchmark measures.

The first is how we summarise performance. A single accuracy figure is misleading here, because
the anomalous events are so rare that a model could score well simply by never flagging one. What
a follow-up team actually needs to know is a pair of numbers. One is \emph{purity}: of the events
the model flags as anomalous, what fraction really are. The other is \emph{completeness}: of the
anomalous events that exist, what fraction the model catches. These are precisely the quantities
tabulated as precision and recall in Table~\ref{tab:perclass}, and their harmonic mean is the F1
score there (macro-F1 being its unweighted average over the six classes). These trade off against each
other---flag more events and you catch more real ones but also let in more false ones---so we fix
the purity to a chosen value, $\bmlPurityTarget$, and report the completeness reached at that
setting as our headline number. This is exactly the form in which a survey team would state a
requirement: ``keep the flagged anomaly sample at least ninety per cent pure, and report how many
detectable anomalies remain.'' This is not the same as planet completeness because
\texttt{NonPSPL} also contains stellar binaries.

The second is a partial correction for how the training set was built. To train efficiently we
subsample the abundant demoted-binary by-products, and each test event carries a weight that
reverses \emph{that} subsampling, applied to purity (which depends on the mix of false positives)
but not to completeness (a per-class property). The weighting undoes one internal sampling step; it
does \emph{not} reweight the deliberately chosen generator class quotas to a survey event-rate
prior.

The third is a reference to judge the anomaly ranking against. For each event we can compute, from
the noiseless light curve and the true parameters, its injected anomaly \dchisq{}---in effect how
detectable the anomaly would be to an observer who already knew the answer. This is a
\emph{truth-informed reference}, not an oracle upper bound: it uses information the survey never
has, and a learned model exploiting other morphology could in principle rank better on real, noisy
data. Comparing BinML against it measures how much of the injected
anomaly ordering the model reconstructs from the raw measurements alone.

\section{Results}\label{sec:results}

We evaluate complete 72-day curves from a held-out pool of $\bmlNpool$ events, split once into a $\bmlNval$-event
validation set, on which the operating threshold is fixed, and a disjoint $\bmlNtest$-event final-test set on
which the results below are computed. On it BinML reaches a completeness of $\bmlCompleteness$ (event-bootstrap
95\% CI $[\bmlCompletenessLo,\bmlCompletenessHi]$) at the fixed anomaly purity of $\bmlPurity$, and an anomaly
average precision of $\bmlAP$. Per-class performance is in Table~\ref{tab:perclass} and the confusion structure in
Figure~\ref{fig:confusion}. The interval covers test-event sampling only, not training-seed or
threshold-calibration uncertainty.

\begin{deluxetable}{lcccc}
\tablecaption{Per-class performance on the independent $\bmlNtest$-event test set. Each event is assigned to the
class with the largest predicted probability (argmax), so these are multi-class figures rather than the
anomaly-detection operating point. Precision, recall and F1 are weighted by $1/\mathrm{keep\_prob}$ to undo the
generation-time subsampling; support is the unweighted number of stored test events with each label.
\label{tab:perclass}}
\tablehead{\colhead{Class} & \colhead{Support} & \colhead{Recall} & \colhead{Precision} & \colhead{F1}}
\startdata
Flat          & \bmlSupFlat          & \bmlFlatR          & \bmlFlatP          & \bmlFlatF \\
PSPL          & \bmlSupPSPL          & \bmlPSPLR          & \bmlPSPLP          & \bmlPSPLF \\
NonPSPL       & \bmlSupNonPSPL       & \bmlNonPSPLR       & \bmlNonPSPLP       & \bmlNonPSPLF \\
PeriodicVar   & \bmlSupPeriodicVar   & \bmlPeriodicVarR   & \bmlPeriodicVarP   & \bmlPeriodicVarF \\
LongPeriodVar & \bmlSupLongPeriodVar & \bmlLongPeriodVarR & \bmlLongPeriodVarP & \bmlLongPeriodVarF \\
Eruptive      & \bmlSupEruptive      & \bmlEruptiveR      & \bmlEruptiveP      & \bmlEruptiveF \\
\hline
\multicolumn{4}{l}{Macro-F1} & \bmlMacroF \\
\enddata
\tablecomments{Of the \texttt{NonPSPL} false positives, $\bmlFpDemotedPct\%$ are events \emph{generated} as
binaries whose anomaly falls below the detectability floor and were therefore labelled \texttt{PSPL}
(\S\ref{sec:detcond}). They remain false positives under the adopted policy; the model's flagging them suggests it responds to
sub-threshold anomaly signal (or to correlated source properties); whether that sensitivity is real needs independent
detection statistics.}
\end{deluxetable}

The largest off-diagonal term is \texttt{NonPSPL}$\rightarrow$\texttt{PSPL} at $\bmlConfNpToPspl\%$ of the true
anomaly row: genuine misses of detectable anomalies, with two sources. Most are wide binaries:
$\bmlNpToPsplWidePct\%$ have separations $s > \bmlSTopBin$, the top bin of the separation prior, which holds only
$\bmlDetWidePct\%$ of the detectable anomalies; these are missed $\bmlWideMissLo$--$\bmlWideMissHi\%$ of the time
at every evidence strength up to $\Delta\chi^2 = 10^6$ and $\bmlWideMissTop\%$ above it, almost all of them at
$q \ge 0.1$ (argmax recall $\bmlRecallWideBigq$, against $\bmlRecallWideSmallq$ for wide binaries at smaller $q$).
At smaller separations the misses concentrate at weak anomalies: those with $\Delta\chi^2 < 2000$ are
$\bmlNotWideWeakDet\%$ of the detectable anomalies there but $\bmlNotWideWeakMiss\%$ of the misses. The reverse
confusion, \texttt{PSPL}$\rightarrow$\texttt{NonPSPL} ($\bmlConfPsplToNp\%$), is what depresses \texttt{NonPSPL}
precision, and $\bmlPsplToNpDemotedPct\%$ of it is demoted binaries. Overall, a detectable binary is missed at a
rate of $\bmlMissedRate$ under six-class argmax (more than half of those misses at the widest separations) and
$\bmlThrMissPct\%$ at the operating threshold; the most dangerous physical confusion, a \texttt{LongPeriodVar} read as
microlensing, occurs at only $\bmlLPVtoPSPL$. The operating threshold is a free parameter, so a follow-up programme
can set the balance between the two error types anywhere along Figure~\ref{fig:pr}.

Because triage uses $P(\texttt{NonPSPL})$ directly, calibration is evaluated for that one-vs-rest probability:
the population-weighted expected calibration error is $\bmlECE$ and the Brier score $\bmlBrier$
(Figure~\ref{fig:calibration}), with no post-hoc scaling. The error is small because $\bmlCalibLowWeightPct\%$ of
the weight sits below 0.1; most of it ($\bmlCalibEceMidPct\%$) comes from scores of 0.1--0.9, where the model is
over-confident (those events are anomalies $\bmlCalibMidFreq$ of the time at a mean score of $\bmlCalibMidScore$,
weighted, and $\bmlCalibMidFreqStored$ at $\bmlCalibMidScoreStored$ among the stored events, so the weighting widens
the gap). Above 0.9, where the threshold sits, the weighted frequency is $\bmlCalibTopFreq$ at a mean score of
$\bmlCalibTopScore$.

\begin{figure*}[t]
\centering
\includegraphics[width=0.88\textwidth]{outputs/figures/calibration.pdf}
\caption{One-vs-rest reliability of $P(\texttt{NonPSPL})$ on the frozen test (population-weighted); marker area
is bin occupancy. Expected calibration error $\bmlECE$, Brier score $\bmlBrier$. The right panel enlarges the
region around the fixed operating threshold.}
\label{fig:calibration}
\end{figure*}

\begin{figure}[tbp]
\centering
\includegraphics[width=0.95\columnwidth]{outputs/figures/confusion.pdf}
\caption{Row-normalised, selection-corrected confusion matrix. The largest off-diagonal,
\texttt{NonPSPL}$\rightarrow$\texttt{PSPL} ($\bmlConfNpToPspl\%$), is genuine recovery loss, mostly on wide
binaries; the smaller \texttt{PSPL}$\rightarrow$\texttt{NonPSPL} term ($\bmlConfPsplToNp\%$) is what depresses
\texttt{NonPSPL} precision.}
\label{fig:confusion}
\end{figure}

Figure~\ref{fig:pr} compares the anomaly ranking with the truth-informed reference of \S\ref{sec:eval}: average
precision $\bmlAPnet$ against the reference's $\bmlAPoracle$. The reference ranks by the same \dchisq{} that
defines the label, so the gap measures how faithfully BinML reconstructs the label-defining quantity from noisy
photometry in a single pass, without fitting a binary-lens model. Argmax recall rises from $\bmlSliceSuzuki$ near
the planet-like mass-ratio break to $\bmlSliceStrongAnom$ for the strongest anomalies.

\begin{figure}[tbp]
\centering
\includegraphics[width=0.85\columnwidth]{outputs/figures/pr_nonpspl.pdf}
\caption{Anomaly (\texttt{NonPSPL}) precision--recall. BinML (solid) tracks the truth-informed \dchisq{}
reference (dashed), computed from the noiseless truth and not fieldable. The marker is the fixed-purity operating
point of the headline number.}
\label{fig:pr}
\end{figure}

The \texttt{NonPSPL} class spans six decades in mass ratio $q$, so Table~\ref{tab:massratio} splits completeness
by regime: $\bmlQStellarComp$ for stellar-binary-like ratios above $10^{-2}$, which dominate the aggregate,
$\bmlQGiantComp$ for $10^{-3}<q<10^{-2}$, $\bmlQNeptuneComp$ for $10^{-4}<q<10^{-3}$ and $\bmlQLowmassComp$ below
$10^{-4}$. These are ratios, not companion masses. Figure~\ref{fig:plane} resolves recovery in the
$(\log s,\log q)$ plane, in the spirit of ground-survey detection-efficiency maps \citep{Tsapras2016}: conditional
on a detectable anomaly, recall has a median of $\bmlEffMedian$ over populated cells, falling to
$\bmlEffMin$--$\bmlEffLowqHi$ in the thinly populated cells at $\log q < -4$ and to
$\bmlEffWideBigqLo$--$\bmlEffWideBigqHi$ in the two well-populated widest-separation cells at $q > 0.1$. The
strong gradient in detectability toward small $q$ is a property of the simulated survey, not of the classifier:
events with $q<10^{-3}$ are recovered at $\bmlSliceDeepPlanet$ (argmax) where they are detectable. Along the individual
axes (Figure~\ref{fig:params}) recall is slightly lower toward faint baselines and weak anomalies; argmax recall remains
$\bmlSliceFaint$ for $m_{\mathrm{base}}>22.5$.

\begin{deluxetable}{lrc}
\tablecaption{Completeness at the fixed operating threshold by lens mass ratio $q$, on the frozen test set. Counts
are anomalies retained by the adopted detectability policy.\label{tab:massratio}}
\tablehead{\colhead{Regime} & \colhead{$N$} & \colhead{Completeness}}
\startdata
Stellar-binary-like, $q>10^{-2}$         & \bmlQStellarN & \bmlQStellarComp \\
High planet-like ratio, $10^{-3}<q<10^{-2}$ & \bmlQGiantN   & \bmlQGiantComp \\
Intermediate ratio, $10^{-4}<q<10^{-3}$  & \bmlQNeptuneN & \bmlQNeptuneComp \\
Low ratio, $q<10^{-4}$                    & \bmlQLowmassN & \bmlQLowmassComp \\
\enddata
\end{deluxetable}

Three noisy-photometry baselines on a common set of $\bmlBaseNtest$ freshly simulated events
(Table~\ref{tab:baselines}) place the network against fieldable alternatives: a fitted single-lens residual
statistic (a \dchisq{}-style score from a Paczynski fit to the noisy F146 curve) and gradient-boosted-tree and
logistic classifiers \citep{scikit-learn} on hand-built morphological features. Every method sees the same events,
the same single band and all 6912 epochs; the comparison set has its own mixture (half of its events are generated
as binaries, nothing is subsampled, $\bmlBasePrev\%$ carry the anomaly label), and the training budgets are
unmatched. BinML scores $\bmlBaseBinml$, against $\bmlBasePspl$ for the fitted residual, $\bmlBaseGbt$ for the trees
and $\bmlBaseLogistic$ for logistic regression.

\begin{deluxetable}{lc}
\tablecaption{Anomaly-detection average precision for BinML and noisy-photometry baselines on a common
$\bmlBaseNtest$-event test set.\label{tab:baselines}}
\tablehead{\colhead{Method} & \colhead{Average precision}}
\startdata
BinML (this work)                       & \bmlBaseBinml \\
Gradient-boosted trees (features)       & \bmlBaseGbt \\
Logistic regression (features)          & \bmlBaseLogistic \\
Fitted single-lens residual ($\dchisq$) & \bmlBasePspl \\
\enddata
\tablecomments{All methods see the same events and the same single band (F146) at full cadence; the feature
classifiers are trained on a disjoint split of the same set. BinML was pretrained on a much larger simulation set,
so this is not a controlled method comparison. Generated by \code{validation/baselines.py}.}
\end{deluxetable}

\begin{figure*}[tbp]
\centering
\includegraphics[width=0.95\textwidth]{outputs/figures/efficiency_plane.pdf}
\caption{Recovery across the $(\log s,\log q)$ plane on the frozen final-test split. \emph{Left:} survey
detectability, the weighted fraction of generated binaries whose anomaly clears the detectability floor.
\emph{Centre:} the classifier's conditional recall given a detectable anomaly. \emph{Right:} $\log_{10}$ of the
effective detectable-event count per cell. Cells with fewer than 30 effective events (left) or effective
detectable events (centre) are blank. The dotted line marks $q=10^{-2}$.}
\label{fig:plane}
\end{figure*}

\begin{figure}[tbp]
\centering
\includegraphics[width=0.98\columnwidth]{outputs/figures/param_dependence.pdf}
\caption{\texttt{NonPSPL} recall versus source baseline magnitude (left) and injected anomaly \dchisq{} (right).
The dotted line at $m_{\rm base}=22.5$ marks the faint-source boundary used in the text.}
\label{fig:params}
\end{figure}

\section{Generalisation and limits}\label{sec:limits}

Held-out numbers can reflect memorised training examples, so we test generalisation directly. The
$\bmlStressN$-million-event stress set is drawn from a disjoint pseudo-random stream (seed base $\bmlSeedStress$
against $\bmlSeedTrain$): a $\bmlStressNatN$-million-event natural-population subset, targeted regimes (enriched
pools that were also used in training, drawn afresh) and out-of-range sweeps that redraw one parameter beyond the
natural priors. The set was scored in July 2026 with the stage-5 checkpoint, the released model's predecessor
(Appendix~\ref{app:infra}); for the released model we regenerated the first shards of each quoted tier with the
set's seeds ($\bmlStressSubN$ events) and scored them with both checkpoints (Table~\ref{tab:stress}). The
regenerated events are a new realisation, not the set's own, because the set's cloud environment was not pinned:
the natural tier labels $\bmlStressNatAnomDeficit\%$ fewer detectable anomalies than the set and the
wide-separation sweep $\bmlStressWidesepAnomDeficit\%$ fewer. Stage~5 on the subset reproduces the set's macro-F1
($\bmlStressFiveSubMacroF$ against $\bmlStressMacroF$) and its other numbers to $\bmlStressSubDiffMed$ at the
median, so we compare checkpoints within the subset. On the natural population the released model reproduces its
held-out performance (macro-F1 $\bmlStressRelMacroF$ against $\bmlMacroF$). Its training is not blind to the
regimes below: the low-$q$ and faint ($m = 23.5$--$25$) regimes were training pools from stage~4, and stage~6
added pools overlapping the wide-separation and sub-day sweeps, so for the released model these numbers measure
how far that coverage carries.

The out-of-range sweeps still expose failures. Anomaly recall falls to $\bmlStressRelWidesep$ for very wide
separations ($s = 5$--$12$; $\bmlStressRelWidesepN$ retained binaries), extending the in-distribution weakness
at the top of the separation prior. Periodic variables with periods comparable to or longer than the season
($60$--$200$~d) are recovered only $\bmlStressRelLongp$ of the time, since within one window they are not
observably periodic. Two failures manufacture spurious anomalies, the mode that matters most for follow-up. Only
$\bmlStressRelShortteSingle$ of detectable sub-day single lenses ($t_{\rm E} = 0.2$--$1$~d;
$\bmlStressRelShortteSingleN$ events) are classified \texttt{PSPL}: $\bmlStressRelShortteSingleNp\%$ are called
anomalies and $\bmlStressRelShortteSingleAlert\%$ cross the frozen threshold
($\bmlStressRelShortteCurSingleAlert\%$ with the peak times drawn for the sub-day timescales, the ``corrected''
rows of the table; the set's per-label value of $\bmlStressSuiteShortte$ mixes in demoted binaries of ordinary
timescale, $\bmlStressShortteDemotedShare\%$ of the weighted \texttt{PSPL} labels in the regenerated tier, which the
model recovers).
Among faint microlensing events ($m = 25$--$27.5$) without a detectable anomaly, the false-anomaly rate rises from
$\bmlStressRelNatMlFa\%$ in the natural population to $\bmlStressRelFaintMlFa\%$ ($\bmlStressRelNatMlAlert\%$ to
$\bmlStressRelFaintMlAlert\%$ above the threshold), and anomaly precision falls to $\bmlStressRelFaintPrec$; faint
photometry, not the sweep's flat-heavy class mix, drives most of that fall (reweighted to the natural class mix
the same events give $\bmlStressFaintPrecNatMix$, against $\bmlStressRelNatNpPrec$). In the low-$q$ regime recall
holds at $\bmlStressRelPlanetNp$ and precision falls to $\bmlStressRelPlanetPrec$ only through prevalence. Each
failure has a plausible physical cause (too few epochs on the peak, a signal that does not complete within the
window, noise comparable to the signal); they are the current model's operating envelope, not established limits of
the method. Events observed without the blue band are still recovered at $\bmlSliceNoBlue$.

\begin{deluxetable*}{llcccc}
\tabletypesize{\footnotesize}
\tablecaption{Stress numbers of the released model and its predecessor. ``Full set'' is the
$\bmlStressN$-million-event stress set scored with the stage-5 checkpoint; ``subset'' is its first shards
regenerated with the same seeds (a new realisation, not the same events) and scored with both checkpoints. Recall
and precision are per label, argmax class, weighted by the inverse keep probability; $N$ counts events with the
class label (for precision, the detectable anomalies; for macro-F1, all events). Single-lens rows keep the generated single lenses with a \texttt{PSPL} label; ``called anomalies'' is
the fraction whose argmax class is \texttt{NonPSPL} (above the frozen threshold in parentheses); ``corrected''
regenerates the sub-day tier with the peak times drawn for the sub-day timescales; ``microlensing'' rows take every
\texttt{PSPL}-labelled event. \code{validation/stress\_rescore\_local.json},
\code{paper/results/stress\_report.json}.\label{tab:stress}}
\tablehead{\colhead{regime} & \colhead{quantity} & \colhead{full set, stage 5 ($N$)} &
\colhead{subset, stage 5} & \colhead{subset, released} & \colhead{$N$ (subset)}}
\startdata
natural & macro-F1 & $\bmlStressMacroF$ ($\bmlStressNatN$M) & $\bmlStressFiveSubMacroF$ & $\bmlStressRelMacroF$ & $\bmlStressSubNat$ \\
natural & anomaly recall & $\bmlStressSuiteNatNp$ ($\bmlStressNatNpN$) & $\bmlStressFiveSubNatNp$ & $\bmlStressRelNatNp$ & $\bmlStressRelNatNpN$ \\
natural & anomaly precision & $\bmlStressSuiteNatNpPrec$ & $\bmlStressFiveSubNatNpPrec$ & $\bmlStressRelNatNpPrec$ & $\bmlStressRelNatNpPrecN$ \\
natural & microlensing: called anomalies & \nodata & \nodata & $\bmlStressRelNatMlFa\%$ ($\bmlStressRelNatMlAlert\%$) & $\bmlStressRelNatPsplN$ \\
low $q$ ($10^{-6}$--$10^{-3}$) & anomaly recall & $\bmlStressSuitePlanetNp$ ($\bmlStressPlanetNpN$) & $\bmlStressFiveSubPlanetNp$ & $\bmlStressRelPlanetNp$ & $\bmlStressRelPlanetNpN$ \\
low $q$ & anomaly precision & $\bmlStressSuitePlanetPrec$ & $\bmlStressFiveSubPlanetPrec$ & $\bmlStressRelPlanetPrec$ & $\bmlStressRelPlanetPrecN$ \\
wide $s$ ($5$--$12$) & anomaly recall & $\bmlStressSuiteWidesep$ ($\bmlStressWidesepN$) & $\bmlStressFiveSubWidesep$ & $\bmlStressRelWidesep$ & $\bmlStressRelWidesepN$ \\
long $P$ ($60$--$200$~d) & \texttt{PeriodicVar} recall & $\bmlStressSuiteLongp$ ($\bmlStressLongpN$) & $\bmlStressFiveSubLongp$ & $\bmlStressRelLongp$ & $\bmlStressRelLongpN$ \\
sub-day $t_{\rm E}$ ($0.2$--$1$~d) & \texttt{PSPL} recall & $\bmlStressSuiteShortte$ ($\bmlStressShortteN$) & $\bmlStressFiveSubShortte$ & $\bmlStressRelShortte$ & $\bmlStressRelShortteN$ \\
 & single lenses: \texttt{PSPL} & \nodata & $\bmlStressFiveSubShortteSingle$ & $\bmlStressRelShortteSingle$ & $\bmlStressRelShortteSingleN$ \\
 & single lenses: called anomalies & \nodata & $\bmlStressFiveSubShortteSingleNp\%$ ($\bmlStressFiveSubShortteSingleAlert\%$) & $\bmlStressRelShortteSingleNp\%$ ($\bmlStressRelShortteSingleAlert\%$) & $\bmlStressRelShortteSingleN$ \\
 & corrected: single lenses \texttt{PSPL} & \nodata & $\bmlStressFiveSubShortteCurSingle$ & $\bmlStressRelShortteCurSingle$ & $\bmlStressRelShortteCurSingleN$ \\
 & corrected: called anomalies & \nodata & $\bmlStressFiveSubShortteCurSingleNp\%$ ($\bmlStressFiveSubShortteCurSingleAlert\%$) & $\bmlStressRelShortteCurSingleNp\%$ ($\bmlStressRelShortteCurSingleAlert\%$) & $\bmlStressRelShortteCurSingleN$ \\
faint ($m = 25$--$27.5$) & \texttt{PSPL} recall & $\bmlStressSuiteFaintPspl$ ($\bmlStressFaintPsplN$) & $\bmlStressFiveSubFaintPspl$ & $\bmlStressRelFaintPspl$ & $\bmlStressRelFaintPsplN$ \\
 & single lenses: \texttt{PSPL} & \nodata & $\bmlStressFiveSubFaintSingle$ & $\bmlStressRelFaintSingle$ & $\bmlStressRelFaintSingleN$ \\
 & microlensing: called anomalies & \nodata & \nodata & $\bmlStressRelFaintMlFa\%$ ($\bmlStressRelFaintMlAlert\%$) & $\bmlStressRelFaintPsplN$ \\
 & anomaly precision & $\bmlStressSuiteFaintPrec$ ($\bmlStressFaintNpN$) & $\bmlStressFiveSubFaintPrec$ & $\bmlStressRelFaintPrec$ & $\bmlStressRelFaintPrecN$ \\
\enddata
\end{deluxetable*}

\paragraph{Cadence.} BinML is a specialist for dense, nearly continuous space-based sampling; archival
ground-survey light curves are outside the evaluated regime. A matched-density test isolates the gap topology of a
ground network such as KMTNet \citep{Kim2016}: at identical visit counts, anomaly recall over $\bmlGapN$ events is
$\bmlGapUniformA$ with uniform sampling against $\bmlGapNightlyA$ through an eight-hour nightly window at
$\bmlGapVisitsA$ visits per day, and $\bmlGapUniformB$ against $\bmlGapNightlyB$ at $\bmlGapVisitsB$. At
$\bmlGapVisitsC$ visits per day the random uniform arm fails too, but a gap-free regular grid with the same visits
still recovers $\bmlGapRegularC$: what makes the shipped model fail outright is empty two-hour bins, of which the
nightly window leaves at least $\bmlGapEmptyNightly\%$. Thinning alone also costs recall (the same events give
$\bmlGapFull$ at full cadence), and at the highest density the regular grid trails the random subsample, so the
arrangement of the visits matters beyond empty bins.
Our own schedule, 15-minute F146 sampling with 6-hour colour visits over one 72-day season, is a legacy
\Rom{}-like design; the current Galactic Bulge Time-Domain Survey specifies F146 every 12 minutes, staggered colour
visits, 66-second exposures and six high-cadence plus four low-cadence seasons \citep{STScI2025GBTDS,ROTAC2025}.
Retraining the whole pipeline at the 12-minute rate makes no difference we can resolve with one seed per arm
(anomaly average precision
$\bmlCadTwelveAp$ against $\bmlCadFifteenAp$ at 15 minutes, no per-class F1 differing by more than
$\bmlCadMaxFOneDelta$). The exposure time is not tested, the colour calibration is tested below
at our exposure time, and the multi-season structure cannot be represented by a single-season input.

\paragraph{The schedule is what the shipped model cannot survive.} The survey as simulated in RMDC26
\citep{RGESPIT2026} pauses F146 for $\bmlGullsPauseHoursLo$--$\bmlGullsPauseHoursHi$~h per
$\bmlGullsSeasonDays$-day season, in six or seven pauses whose phases differ between seasons. Essentially none of our
training events contains an empty mid-season F146 bin, and imposing the measured pauses on our own held-out seasons, with
nothing else changed, collapses the shipped model (Table~\ref{tab:gap}): single-lens recall falls from
$\bmlGapShippedPsplClean$ to $\bmlGapShippedPsplSched$, mostly to \texttt{NonPSPL}, and \texttt{Flat} recall from
$\bmlGapShippedFlatClean$ to $\bmlGapShippedFlatSched$; a single gap of 1--2~h at mid-season already lowers
single-lens recall from $\bmlGapSensNone$ to $\bmlGapSensTwo$. The remedy is an augmentation: blanking random
runs of 1--12~h during a 12-epoch warm-start fine-tune raises held-out average precision under the pauses from
$\bmlGapApShipped$ to $\bmlGapApRandNorelabel$ at a clean-data cost of $\bmlGapApCleanShipped$ to
$\bmlGapApCleanRandNorelabel$ (with the caustic-in-gap relabel off, as in the measured-season arm),
and training on one of RMDC26's six measured seasons per example gives $\bmlGapApSeasons$; the two are within the
spread of three training seeds ($\bmlSeedHeldApRange$) on our seasons, while on RMDC26 the measured pauses lead
in recall (\S\ref{sec:gulls}). That relabel rule relabels a binary as
\texttt{PSPL} when the bins around its onset are blanked, but applied to the epochs that survive the pauses it
gives $\bmlTruthSeasonsOn\%$ of binary presentations a single-lens label although their surviving anomaly would
still be claimed, whereas with it off $\bmlTruthSeasonsOff\%$ keep a binary label the surviving epochs do not
support. The shipped checkpoint's input contract is therefore \emph{continuous F146}; for a gapped schedule use
the recommended checkpoint of \S\ref{sec:gulls}.

\begin{deluxetable*}{lcccccc}
\tablecaption{Our own held-out seasons ($\bmlGapHeldN$ scored events) with the measured RMDC26 pauses imposed
(event $i$ receives the empty bins of high-cadence season $i \bmod 6$); ``clean'' is the same set without pauses.
Rows other than the shipped reference and g08e12 are 12-epoch warm-start fine-tunes of the shipped weights on one
natural-prior pool; one seed each. Generated from \code{validation/gulls/schedule\_finetune.json}.\label{tab:gap}}
\tablehead{\colhead{augmentation} & \colhead{AP clean} & \colhead{AP pauses} & \colhead{PSPL\tablenotemark{a}} &
           \colhead{Flat} & \colhead{NonPSPL} & \colhead{macro-F1}}
\startdata
shipped & 0.959 & 0.169 & 0.03 & 0.17 & 0.84 & 0.375 \\
gap augmentation, g08e12 (other pool) & 0.952 & 0.896 & 0.89 & 0.97 & 0.92 & 0.875 \\
random gaps & 0.951 & 0.910 & 0.91 & 0.98 & 0.92 & 0.896 \\
random gaps, relabel off & 0.952 & 0.911 & 0.91 & 0.97 & 0.92 & 0.899 \\
first-season mask, relabel off & 0.952 & 0.568 & 0.41 & 0.76 & 0.98 & 0.674 \\
measured seasons, relabel off & 0.948 & 0.921 & 0.91 & 0.97 & 0.94 & 0.899 \\

\enddata
\tablenotetext{a}{PSPL, Flat and NonPSPL are argmax recalls under the pauses. The first-season arm trained on one
season's pauses and matches one RMDC26 season in six.}
\end{deluxetable*}

\paragraph{Prevalence.} Purity depends on the prevalence of anomalies in the input stream, which our mixture sets
at $\bmlPriorPiZero\%$ and a real survey stream sets far lower. Under a pure prior shift, completeness at the
fixed threshold stays at $\bmlCompleteness$ while purity falls from $\bmlPurity$ to $\bmlPurityOnePct$ at a 1\%
prevalence and $\bmlPurityTenthPct$ at 0.1\% (Figure~\ref{fig:prior}); re-tuning the threshold on the validation
rows to hold $\bmlPrevPurTarget$ purity instead leaves completeness $\bmlPrevCompOne$ at 1\% and
$\bmlPrevCompTenth$ at 0.1\%. A deployment should plan against the figure matching its own prevalence, and this
transfer assumes events that look like ours in different proportions, not a different survey.

\begin{figure*}[tbp]
\centering
\includegraphics[width=0.86\textwidth]{outputs/figures/prior_sensitivity.pdf}
\caption{Prior sensitivity of the operating point. At the fixed threshold, completeness (dotted) is independent of
anomaly prevalence, but purity (solid) and the alert burden (dash-dot, right axis) are not. The right panel shows
held-out completeness when the threshold is re-selected on the validation split for each assumed prevalence at
purity $\bmlPrevPurTarget$. The dashed line marks our mixture's $\bmlPriorPiZero\%$ prevalence.}
\label{fig:prior}
\end{figure*}

\paragraph{Colour calibration and the detectability floor.} An audit against the current Roman tables
\citep{RomanTechInfo} finds our F146 zeropoint accurate to $0.03$~mag but the F087 and F213 zeropoints optimistic by
about $0.10$ and $0.14$~mag, the F087 saturation limit too faint (at equal exposure F087 saturates about
$1.3$~mag \emph{brighter} than F146, not fainter), and the thermal backgrounds not reproduced. Regenerating two
held-out shards ($\bmlRefN$ events) with the audited colour bands and scoring both versions with the shipped
model, the anomaly channel barely moves (completeness $\bmlRefColCompAudited$ against $\bmlRefColCompTrained$,
average precision $\bmlRefColApAudited$ against $\bmlRefColApTrained$) while the variable classes lose precision,
not recall: the model calls more flat sources periodic and more single lenses eruptive (periodic-variable F1
$\bmlRefColFonePerTrained$ to $\bmlRefColFonePerAudited$). Short fine-tunes on each calibration show that this
flat-versus-periodic boundary follows the calibration a model is trained on, by an amount that depends on where
training stops; the anomaly channel does not. The
corrected values are recorded in the released code. The class boundaries also depend on the minimum detectable
amplitude, fixed at $0.02$~mag as a stand-in for correlated photometric systematics. Regenerating the same shards
with the floor at $0.01$ and $0.05$~mag (partly overlapping draws: $\bmlRefFloorShareLo\%$ and $\bmlRefFloorShareHi\%$
of the adopted-floor events recur) and scoring them with the shipped model at its frozen threshold, halving
the floor raises the anomaly prevalence from $\bmlRefFloorPrevMid\%$ to $\bmlRefFloorPrevLo\%$ and lowers
completeness from $\bmlRefFloorCompMid$ to $\bmlRefFloorCompLo$ while purity rises from $\bmlRefFloorPurMid$ to
$\bmlRefFloorPurLo$, and raising it moves all three the other way, with average precision staying between
$\bmlRefFloorApMin$ and $\bmlRefFloorApMax$: a different floor leaves the anomaly ranking largely intact but moves
the operating point. The completeness and purity we quote are conditional on the $0.02$~mag floor; a survey
requirement should rest on a model retrained at the floor its photometry supports.

\section{Cross-simulator validation on an independent \Rom{} simulation}\label{sec:gulls}

Every number so far comes from our own simulator. RMDC26 \citep{RGESPIT2026} simulates the survey with GULLS
\citep{Penny2013,Penny2019} and the SynthPop Galactic model \citep{Kluter2025}, with its own schedule, photometry
and population. No RMDC26 event entered training and no deployed threshold was chosen on it, but it was used to
diagnose the residual false alarms, to set the upper end of the single-lens source-size prior and to choose among
$\bmlGullsNcheckpoints$ fine-tuned checkpoints, so the recommended checkpoint's numbers are the best of several on
this set and are optimistic.

\subsection{Data, selection, and input contract}

RMDC26 contains $\bmlGullsNcat$ events: $\bmlGullsNcatSingle$ single lenses (1S1L), $\bmlGullsNcatPlanet$
planetary lenses (1S2L) and $\bmlGullsNcatPlanetBin$ planetary lenses with a second source (2S2L), the planetary
classes pre-selected by GULLS' own detection statistic. It has no pure binary-source class, so the
\citet{Gaudi1998} degeneracy remains untested, and the source is flagged binary in $\bmlGullsBinSrcPct\%$ of scored
2S2L events, so a \texttt{NonPSPL} call on a 2S2L event may respond to its second source rather than its planet; and its anomalies are planetary (median mass ratio
$\bmlGullsQmedPlanet$; about $\bmlCgRmStellarFrac\%$ have $q > 10^{-2}$) whereas $\bmlDetStellarPct\%$ of our
test set's detectable anomalies are stellar-mass-ratio binaries, so transfer recall is not comparable to the
completeness of \S\ref{sec:results}. Selection uses catalogue parameters only and is applied identically to the
three classes: a host single-lens peak amplitude of at least $0.1$~mag and $1 \le t_{\rm E} \le 300$~d, our
natural prior's support ($\bmlGullsSubdaySingle\%$ of RMDC26 single lenses are sub-day). That leaves
$\bmlGullsEligibleN$ events, of which $\bmlGullsScoredPct\%$ ($\bmlGullsN$: $\bmlGullsNSingle$ 1S1L,
$\bmlGullsNPlanet$ 1S2L, $\bmlGullsNPlanetBin$ 2S2L) peak inside one of the six high-cadence seasons and are
scored; $\bmlGullsBetweenPct\%$ peak between seasons (\S\ref{sec:gulls:seasons}) and $\bmlGullsLowcadPct\%$ in a
low-cadence season, outside our support. The amplitude cut is on the host, so the scored planetary events are
dominated by resonant and central-caustic geometries ($\bmlGullsWidePass\%$ of 1S2L events with $s > 2$ pass,
against $\bmlGullsResonantPass\%$ at $0.5 \le s < 2$). Some inputs lie outside our training priors: baselines
brighter than our 20~mag limit ($\bmlGullsBrightSingle\%$ of single lenses), blend fractions below our prior's
minimum, and much shorter single lenses ($\bmlGullsShortSingle\%$ have $t_{\rm E} < 3$~d,
$\bmlGullsShortSingleW\%$ weighted by event rate, against $\bmlPriorShortPct\%$ of our prior). The model is given
the 72-day window from the season boundary in all three bands, with the baseline taken as the median F146
magnitude far from the catalogue peak (a real-time pipeline would have neither), and observations pooled to one
value per 15-minute epoch as in the training cache; binning by raw observation instead moves the shipped
checkpoint's single-lens false-alarm rate from $\bmlGullsFaShipped\%$ to $\bmlGullsRawpoolFa\%$. RMDC26's colour
visits also displace one of eight F146 epochs in about $\bmlGullsPartialBinsPct\%$ of the 2-h bins, an occupancy
our training data never contain; setting it to full on $\bmlOccN$ real inputs lowers the recommended checkpoint's
single-lens false alarms from $\bmlOccFaObs\%$ to $\bmlOccFaOne\%$ and its planetary recall from $\bmlOccRecObs$
to $\bmlOccRecOne$.

\subsection{Results}\label{sec:gulls:results}

Table~\ref{tab:gulls} gives, on the identical $\bmlGullsN$ events for every checkpoint, the single-lens
false-alarm rate at the frozen threshold and planetary recall at matched single-lens false-alarm budgets, the
comparison a threshold shift cannot fake. Rates are per simulated event and, where marked, weighted by RMDC26's
event-rate weight; the unweighted sample over-represents short single lenses, and false alarms rise with
timescale, so the two differ, and we say where a conclusion depends on the weighting. Intervals are event
bootstraps and cover the evaluation sample, not training variation.

The shipped checkpoint flags $\bmlGullsFaShipped\%$ of single lenses, its first choice across single lenses split
between \texttt{NonPSPL} and \texttt{PeriodicVar}, the signature of the empty mid-season tokens
(\S\ref{sec:limits}). Gap augmentation brings the rate to $\bmlGullsFaGapaware\%$. Its residual false alarms rise
with $\rho/|u_0|$, the ratio of source radius to impact parameter, from $\bmlGullsRhoLoGapaware$ for point-like
sources to $\bmlGullsRhoHiGapaware$ above 3: these are high-magnification events, and the likely cause is that our
generator's single lenses were point sources, so a finite-source single lens resembled no training \texttt{PSPL}. Adding
finite-source single lenses (\textsc{VBBinaryLensing} \citep{Bozza2018}, $\rho$ log-uniform up to 5) and
fine-tuning for 12 further epochs from the gap-augmented checkpoint lowers the rate to $\bmlGullsFaFspl\%$ and the
largest-ratio bin to $\bmlGullsRhoHiFspl$, and raises planetary recall at the $\bmlGullsBudgetMid\%$ budget from
$\bmlGullsRecGapaware$ to $\bmlGullsRecFspl$. A control with the same pool shape, recipe and epochs but
point-source single lenses separates the physics from the extra training: the training alone reaches
$\bmlGullsFaCtrl\%$ and $\bmlGullsRecCtrl$, so in false alarms and at low budgets most of the gain is the extra
training, while at higher budgets and in mean recall it is the physics, which also lowers false alarms further in the
two largest $\rho/|u_0|$ bins (Table~\ref{tab:rho}; per simulated event; weighted by event rate the physics resolves
nothing, and the largest-ratio bin stays near $\bmlGullsRhoHiFspl$, so the missing physics explains only part of the
residual). Training on
the measured per-season pauses instead of random gaps also transfers: on one pool it
lowers false alarms from $\bmlGullsFaRandNorelabel\%$ to $\bmlGullsFaSeasons\%$ and raises recall at the budget
from $\bmlGullsRecRandNorelabel$ to $\bmlGullsRecSeasons$, leading in $\bmlGullsSeasonWins$ of the
$\bmlGullsNseasons$ seasons. Matching GULLS' noise, which reaches $\bmlGullsNoiseFaint$ times ours at
24--25~mag, gave a checkpoint that was worse at every matched budget ($\bmlGullsFaNoisy\%$; recall
$\bmlGullsRecNoisy$).

The recommended checkpoint, marked (rec.) in Table~\ref{tab:gulls}, combines the finite-source recipe with the
measured pauses (and a smooth finite-source magnification, which changes nothing measurable). Per simulated event
its planetary recall is within $\bmlGullsCombDiffMax$ of the finite-source checkpoint's at every budget and it
flags slightly more single lenses at the frozen threshold ($\bmlGullsFaCombined\%$ against $\bmlGullsFaFspl\%$);
weighted by event rate it is ahead at every budget, by $\bmlDiffCombRecWMin$--$\bmlDiffCombRecWMax$ for 1S2L, and
it is better on our own held-out seasons with the pauses imposed (average precision $\bmlSidecarHeldAp$ against
$\bmlFsplHeldAp$). Two further seeds of the recipe, trained after the choice (the last rows of the table), trail
the finite-source checkpoint per event and tie it weighted, so RMDC26 does not separate the two recipes; the three
seeds give the scale of the training noise behind every comparison here (Table~\ref{tab:seeds},
Appendix~\ref{app:rmdc26}): planetary recall
$\bmlSeedRecLo$--$\bmlSeedRecHi$ at the $\bmlGullsBudgetMid\%$ budget and false alarms
$\bmlSeedFaLo$--$\bmlSeedFaHi\%$. We read a difference between single runs as unresolved when it is smaller than
that range at the same budget and weighting. On that reading the measured pauses lead random gaps in recall in both
weightings (in false alarms per simulated event only), the extra training carries the false-alarm gain, the finite-source physics carries the high-budget
recall gain per event only, and the recommended recipe's lead over the finite-source one is not established.
Removing the colour bands lowers planetary recall at matched budget by
$\bmlGullsColourGainLo$--$\bmlGullsColourGainHi$ for each of the $\bmlGullsNcolour$ checkpoints ablated, but for the
recommended one most of the gain is on events without a claimable anomaly (flag rate $\bmlGullsColUndetOne$ to $\bmlGullsColUndetThree$, against recall $\bmlGullsColDetOne$ to
$\bmlGullsColDetThree$ on those with one), so this is not evidence that the colour bands carry anomaly
information across simulators.


\begin{deluxetable*}{lcccccc}
\tablecaption{Cross-simulator transfer on $\bmlGullsN$ RMDC26 events ($\bmlGullsNSingle$ 1S1L, $\bmlGullsNPlanet$
1S2L, $\bmlGullsNPlanetBin$ 2S2L), identical inputs for every row. FA: percentage of single lenses over the frozen
threshold $\bmlThreshold$. R$_x$: 1S2L recall at the threshold that yields a single-lens false-alarm rate of $x\%$
on this population (the $\bmlGullsBudgetMid\%$ and $\bmlGullsBudgetHi\%$ budgets are close to the frozen-threshold
rates of the first finite-source round, $\bmlGullsFaFsplOne\%$, and of gap augmentation, $\bmlGullsFaGapaware\%$,
chosen after seeing the results). Columns marked (w) weight each event by its RMDC26 event
rate (effective sample sizes $\bmlGullsNeffSingle$ single lenses, $\bmlGullsNeffPlanet$ 1S2L). Rows below gap
augmentation add 12 epochs from it with point-source single lenses (control), finite-source single lenses, and
finite sources plus the measured pauses (recommended); the last two rows repeat the recommended recipe with two
other seeds, trained after the choice. One seed each otherwise. The rows up to the finite-source one were trained
with the approximate caustic-in-gap relabel on (\S\ref{sec:limits}); the last three have it off. The false-alarm rate by source size and the
seed spread are in Appendix~\ref{app:rmdc26}. Generated from
\code{validation/gulls/transfer\_tradeoff\_all.json}.\label{tab:gulls}}
\tablehead{\colhead{checkpoint} & \colhead{FA} & \colhead{FA (w)} & \colhead{R$_{\bmlGullsBudgetLo}$} &
           \colhead{R$_{\bmlGullsBudgetMid}$} & \colhead{R$_{\bmlGullsBudgetHi}$} & \colhead{R$_{\bmlGullsBudgetMid}$ (w)}}
\startdata
shipped & 35.6 & 38.3 & 0.122 & 0.214 & 0.313 & 0.144 \\
gap augmentation (g08e12) & 11.6 & 13.5 & 0.185 & 0.295 & 0.457 & 0.336 \\
\quad + point-source control & 6.4 & 8.6 & 0.244 & 0.361 & 0.485 & 0.374 \\
\quad + finite source & 4.8 & 7.6 & 0.265 & 0.390 & 0.533 & 0.383 \\
\quad + finite source, pauses (rec.) & 6.3 & 8.4 & 0.273 & 0.391 & 0.531 & 0.407 \\
\qquad same recipe, seed 2 & 4.7 & 5.8 & 0.232 & 0.349 & 0.503 & 0.386 \\
\qquad same recipe, seed 3 & 7.3 & 9.7 & 0.233 & 0.358 & 0.514 & 0.385 \\
% achieved single-lens false-alarm rate within 0.05 points of each budget

\enddata
\end{deluxetable*}

\paragraph{Operating point.} The frozen threshold was chosen on gap-free, point-source data. For the recommended
checkpoint we set the threshold at the 90\% purity target on our own seasons with the measured pauses imposed,
without reading RMDC26, using the whole finite-source held-out pool because a single 20\% slice scatters the
threshold enough to move RMDC26's false-alarm rate between $\bmlSidecarFaLo$ and $\bmlSidecarFaHi\%$ (68\% of 200
random slices). At the
resulting $\bmlSidecarThr$ the checkpoint flags $\bmlSidecarFa\%$ of RMDC26 single lenses ($\bmlSidecarFaW\%$
weighted) with planetary recall $\bmlSidecarRecPlanet$ / $\bmlSidecarRecPlanetBin$ (1S2L / 2S2L); the three seeds
span $\bmlSeedOwnFaLo$--$\bmlSeedOwnFaHi\%$ and $\bmlSeedOwnRecLo$--$\bmlSeedOwnRecHi$.

\subsection{What ``recall'' means across two label ontologies}\label{sec:gulls:labels}

GULLS assigns classes by generator; BinML's are observational (\S\ref{sec:detcond}). Applying our labelling rule
to RMDC26's noise-free light curves, with GULLS' errors as $\sigma$, to the first $\bmlGullsTruthNSingle$ scored
single lenses and $\bmlGullsTruthNPlanet$ events of each planetary class, $\bmlGullsUndetPlanet\%$ of 1S2L and
$\bmlGullsUndetPlanetBin\%$ of 2S2L events carry no anomaly the policy would claim within one season
($\bmlGullsFailChi\%$ fail the $\Delta\chi^2$ cut, $\bmlGullsFloorVeto\%$ pass it but fall under the floor). This
depends on the floor, the window and our selection; it is not a statement about what a survey could detect. The
recommended checkpoint flags $\bmlGullsFlagUndet\%$ of the planetary events without a detectable anomaly; on those
with one its recall at the frozen threshold is $\bmlGullsRecDetPlanet$ / $\bmlGullsRecDetPlanetBin$, against
$\bmlGullsRecGenPlanet$ / $\bmlGullsRecGenPlanetBin$ on generator labels, rising from $\bmlGullsRecChiLo$ for the
weakest claimable anomalies to $\bmlGullsRecChiHi$ above $\Delta\chi^2 = 10^5$. It also flags
$\bmlGullsFlagFloorVeto\%$ of the planetary events whose anomaly is significant but below the $0.02$~mag floor,
at a similar rate however far below it, and $\bmlGullsSingleMisfitK$ of the $\bmlGullsSingleMisfitN$ single lenses
whose static refit shows a significant sub-floor misfit, so the response to sub-floor structure is not specific
to planets. Lowering the floor to $0.005$~mag raises the fraction of 1S2L events with a claimable anomaly from
$\bmlGullsDetPlanet\%$ to $\bmlGullsFloorDetLo\%$; raising it to $0.05$~mag lowers it to $\bmlGullsFloorDetHi\%$.

\subsection{Events whose peak falls between seasons}\label{sec:gulls:seasons}

The single-season contract leaves out the $\bmlGullsBetweenPct\%$ of eligible events that peak between seasons.
For the first $\bmlMsN$ such events of each class we scored the high-cadence season on each side of the gap and
kept the larger anomaly probability. A high-cadence season borders the gap for $\bmlMsScorablePlanet\%$ of the
between-season 1S2L events, and a claimable anomaly is visible in a wing for $\bmlMsDetPlanet\%$ of those,
against $\bmlGullsDetPlanet\%$ in-season, so the unobserved gap costs $\bmlMsDetLossPlanet$ points of claimable
anomalies (2S2L: $\bmlMsDetPlanetBin\%$, $\bmlMsDetLossPlanetBin$ points). At the frozen threshold the recommended
checkpoint flags $\bmlMsFa\%$ of these single lenses and recalls $\bmlMsRecDetPlanet$ / $\bmlMsRecDetPlanetBin$
of the planetary events with a claimable anomaly in a wing, against $\bmlGullsRecDetPlanet$ /
$\bmlGullsRecDetPlanetBin$ in-season; at a matched $\bmlGullsBudgetMid\%$ single-lens budget the recall is
comparable ($\bmlMsRecBudgetPlanet$ against $\bmlGullsRecDetBudgetPlanet$, a threshold set on only
$\bmlMsBudgetExceed$ exceedances), so scoring seasons around a gap would
need its own operating point. A whole-mission model would add these events and the low-cadence seasons; no model
recovers anomalies that fall in the unobserved gaps.

\subsection{The partial-season cascade on RMDC26}\label{sec:gulls:cascade}

We repeated the primary protocol of \S\ref{sec:cascade} on RMDC26 with the recommended checkpoint at the frozen
threshold: the first $\bmlCgNsingle$ scored single lenses and $\bmlCgNplanet$ of each planetary class revealed in
half-day steps and scored from F146 alone, with the onset taken from our labelling rule on GULLS' noise-free
curve. Of the $\bmlCgElig$ eligible binaries, $\bmlCgPrem\%$ are flagged before the onset (95\% interval
$\bmlCgPremLo$--$\bmlCgPremHi\%$), $\bmlCgDet\%$ by the end of the season, and the median lag is $\bmlCgLag$~d.
Compared stratum by stratum with the in-house scan of \S\ref{sec:cascade} (shipped model, our simulator, whose
anomalies are mostly stellar), premature alerts are about as frequent at high planet-like
($10^{-3} < q \le 10^{-2}$) and intermediate ($10^{-4} < q \le 10^{-3}$) mass ratios (Fisher
$p \ge \bmlCgPremFisherMin$), detection is lower in both strata ($\bmlCgRmGiantDet\%$ against
$\bmlCgInGiantDet\%$, nominally resolved at $p = \bmlCgGiantDetFisher$, one of six comparisons here; $\bmlCgRmNeptuneDet\%$ against $\bmlCgInNeptuneDet\%$ on
$\bmlCgInNeptuneN$ in-house events, unresolved), and alerts come later by $\bmlCgNeptuneLagDiff$~d at
intermediate ratios (bootstrap 95\% interval $\bmlCgNeptuneLagDiffLo$--$\bmlCgNeptuneLagDiffHi$~d, at the edge of
resolution) and
$\bmlCgGiantLagDiff$~d at high ones (unresolved); a different checkpoint on a different simulator, so the
differences mix the two. $\bmlCgRmLowmassFrac\%$ of RMDC26's eligible anomalies have $q \le 10^{-4}$, where
detection falls to $\bmlCgRmLowmassDet\%$. Scanning single lenses gives the alert burden: $\bmlCgSingleAlert\%$
($\bmlCgSingleAlertW\%$ weighted) raise an alert at some point in the season, $\bmlCgSingleAlertDay$ per 1,000
single lenses per day (slightly optimistic: on the complete season this sample flags $\bmlCgSingleFullWindow\%$ of
single lenses with F146 alone, against $\bmlCgSingleAllFone\%$ for all scored ones), several times the rate on our own simulator with the same checkpoint
($\bmlRefRecSinglePct\%$). If 1\% of events were planetary binaries, $\bmlCgPurOne\%$ of alerts would come from a
claimable anomaly ($\bmlCgPurFive\%$ at 5\%; $\bmlCgPurRmdc\%$ at RMDC26's own rate-weighted planetary fraction);
RMDC26 contains no variables or flat sources, so a real stream would be less pure. At the recalibrated threshold
single-lens alerts fall to $\bmlCgSingleAlertCal\%$ and purity at 1\% rises to $\bmlCgPurOneCal\%$, at the cost
of detection ($\bmlCgDetCal\%$) and lag ($\bmlCgLagCal$~d). The cascade's low premature-alert rate holds on
another simulator, but at realistic prevalence its alerts need a second stage that removes single-lens
contaminants.

\subsection{What this test does and does not establish}

It establishes that the shipped model's failure on a gapped schedule is a property of the schedule, reproducible
without GULLS data and reduced by augmentation, preferably with the measured pauses; that continued training on a
pool rich in high-magnification single lenses accounts for most of the further reduction in false alarms; and
that finite-source single lenses, a physics our generator lacked, lower false alarms further at large
$\rho/|u_0|$ and raise recall at high budgets, per simulated event. It does not establish a \Rom{} planet yield:
RMDC26's population, detection cut and our selection all differ from our training population, neither simulator
contains real systematics, every fine-tune but the recommended recipe is a single seed, and RMDC26 was used to
choose among checkpoints. Sub-day single lenses, outside our natural timescale prior, are flagged by the
recommended checkpoint less often than in-support ones ($\bmlGullsSubdayFa\%$ of $\bmlGullsSubdayN$ against
$\bmlGullsFaCombined\%$), their most probable class being mostly \texttt{PeriodicVar}, and at $0.25$--$1$~d the released model sends
$\bmlGullsSubdayShippedMatchedFa\%$ of RMDC26's single lenses above its threshold against
$\bmlStressShortteCurMatchedAlert\%$ in our own sweep at the same timescales, so timescale alone does not explain
that in-house failure; a model for free-floating planets needs sub-day timescales in its prior.

\section{Discussion}\label{sec:discussion}

BinML is intended as a cheap first-pass triage layer before binary-lens modelling: the completeness-at-fixed-purity
framing exposes the trade-off a survey team must set, and the forward pass is fast enough for repeated
survey-scale scoring. The evaluation uses synthetic photometry, a simulator-provided baseline, a truth-informed
onset and half-day updates, so this is a computational benchmark rather than a deployable broker. A median alert
$\bmlCascMedianLagNonPrem$~days after onset arrives once a short caustic crossing has finished, so for brief
planetary anomalies BinML is a post-hoc identifier rather than a trigger for the feature itself; a few-day latency
still serves anomalies that are not brief (extended or repeating caustic structure, long-timescale binaries) and
everything that depends on recognising the event rather than the feature: scheduling denser photometry for the
rest of an ongoing event, initiating colour or parallax campaigns, and prioritising which light curves receive
expensive modelling. Triggering on the anomaly itself would require evaluating at the native cadence, which our
grid sensitivity (\S\ref{sec:cascade}) does not cover. BinML's input is one contiguous season, which leaves out
the events that peak between seasons ($\bmlGullsBetweenPlanetLo$--$\bmlGullsBetweenPlanetHi\%$ of RMDC26's
eligible planetary events; \S\ref{sec:gulls:seasons}) and in low-cadence seasons; a whole-mission model taking the full
light curve, with a mission-length simulation and a hierarchical input, is the natural next version.

\subsection{Ablating the two mechanisms}\label{sec:ablations}

Both methodological claims are testable by removing the mechanism and retraining. Each ablation trains arms from
scratch on identical shards with identical hyper-parameters and seed, differing only in the stated variable, and
evaluates them on a common held-out set from an unseen seed base; these are small-scale runs (16 shards,
$\bmlAbLabelEpochs$ epochs), so absolute values sit below the headline numbers and the comparison between arms is
the result.

\paragraph{Detectability-conditioned labelling.} With the training label the only intervention (both arms use the
same policy-conditioned class weights), macro-F1 against the policy-conditioned label falls from
$\bmlAbLabelObsMacroFOne$ to $\bmlAbLabelGenMacroFOne$ when the network is trained on generator intent, a gap of
$\bmlAbLabelGap$ (95\% CI $[\bmlAbLabelGapLo,\bmlAbLabelGapHi]$, paired bootstrap over $\bmlAbLabelNeval$
events); against the generator label the ordering reverses ($\bmlAbLabelObsMacroFOneGen$ versus
$\bmlAbLabelGenMacroFOneGen$), so each arm wins on the ontology it was trained for. The damage concentrates where
the relabelling acts (\texttt{PSPL} F1 $\bmlAbLabelObsPsplFOne$ to $\bmlAbLabelGenPsplFOne$, \texttt{NonPSPL}
$\bmlAbLabelObsNonpsplFOne$ to $\bmlAbLabelGenNonpsplFOne$), and the generator-intent arm behaves differently: at
the frozen threshold it reaches completeness $\bmlAbLabelGenComp$ at a purity of $\bmlAbLabelGenPur$, because it
was trained to call anomalous the sub-threshold binaries that make up $\bmlGenNpDemoted\%$ of the stored events
of the final-test split. Both arms are single training seeds.

\paragraph{The truncation augmentation: unresolved.} At a single threshold the second ablation does not identify
a timing benefit. Counted against the generator's $7.2$-day onset grid on $N=\bmlAbCascN$ binaries, removing the
augmentation moves the premature rate in opposite directions depending on the rule: at the frozen threshold the
augmented model is better ($\bmlAbCascOnPremThr$ versus $\bmlAbCascOffPremThr$), under argmax worse
($\bmlAbCascOnPremArgmax$ versus $\bmlAbCascOffPremArgmax$). But the grid records only that the anomaly was
undetectable at the previous cut, and almost all of the augmented arm's premature alerts fall inside that window
(alerts before it number $\bmlAbCascOnUnambThr$ against $\bmlAbCascOffUnambThr$ at the threshold and
$\bmlAbCascOnUnambArgmax$ against $\bmlAbCascOffUnambArgmax$ under argmax, although the first-detectable onset lies
earlier still for $\bmlAbFineBeforePrevCutPct\%$ of the in-house binaries), so the argmax reversal is not
established. The augmentation is also expensive: detection at the same threshold falls from $\bmlAbCascOffDet$ to
$\bmlAbCascOnDet$ and static macro-F1 from $\bmlAbCascOffMacroFOne$ to $\bmlAbCascOnMacroFOne$, consistent with
the augmented model saying \texttt{NonPSPL} less often. An exploratory risk--coverage comparison on stored traces
(Figure~\ref{fig:cascade_matched}), thresholding each arm to the same detection fraction on the same $\bmlMatchN$
events, gives the augmented arm fewer premature alerts from $0.50$ through $0.90$ detection ($\bmlMatchOnSixty\%$
against $\bmlMatchOffSixty\%$ at $0.60$), reversing at $0.95$; because thresholds and outcomes use the same events
and one seed per arm, it is descriptive, not confirmatory. A causal timing claim needs thresholds selected on
disjoint prefixes, held-out paired evaluation and repeated seeds.

\begin{figure}[tbp]
\centering
\includegraphics[width=0.98\columnwidth]{outputs/figures/cascade_matched.pdf}
\caption{Exploratory risk--coverage comparison between the truncation-augmented and unaugmented arms, thresholded
to the same detection count at each point; bars are Wilson 95\% intervals on the eligible-event denominator.
Thresholds and outcomes use the same $N=\bmlMatchN$ events.}
\label{fig:cascade_matched}
\end{figure}

The natural next step is validation on real events in the model's own regime, with faint-source false anomalies,
the baseline estimate and the colour calibration as the priorities.

\section{Conclusion}\label{sec:conclusion}

BinML is a compact multi-band classifier for six classes of synthetic, progressively revealed \Rom{}-like light
curves. On the adopted class mixture it reaches completeness $\bmlCompleteness$ at purity $\bmlPurity$ for
detectable \texttt{NonPSPL} events, values conditional on the simulation and its priors rather than planet
yields. Revealed at half-day resolution in F146, $\bmlCascDetFrac\%$ of eligible simulated binaries alert,
$\bmlCascPrematurePct\%$ of the eligible binaries before a truth-informed onset, with a median delay of
$\bmlCascMedianLagNonPrem$~days; with all three bands a mixture of every class raises $\bmlRefStreamAlertsDay$ alerts
per 1,000 events per day, $\bmlRefStreamPurity\%$ of them detectable anomalies at the simulated prevalence
($\bmlRefStreamFPurity\%$ with F146 alone). On an independent
simulation of the survey the shipped model fails, mostly because of the survey's F146 observing pauses, and a
fine-tune trained on those pauses reduces the failure (\S\ref{sec:gulls}). The remaining tests are retraining under
another detectability floor, an estimated-baseline ablation, repeated training seeds for the shipped model, a
streaming threshold calibrated on disjoint prefixes, a whole-mission model, the binary-source degeneracy, sub-day
timescales in the prior, and validation on real photometry. The released code and weights provide a rebuildable
benchmark.

\begin{acknowledgments}
I thank the gravitational lensing group at the Astronomisches Rechen-Institut (ARI), Zentrum für
Astronomie der Universität Heidelberg, and in particular Yiannis Tsapras, for scientific feedback
and discussions that shaped this work.
Simulations were run on Amazon Web Services (AWS) EC2, and
model development used a combination of local Apple Silicon hardware and cloud GPU credits from
AWS and Modal; I gratefully acknowledge both for computational credits. This work made use of
\textsc{VBBinaryLensing} \citep{Bozza2018} and the OGLE variable-star catalogues
\citep{Soszynski2013,Soszynski2014}. In the interest of transparency, I note that AI assistants
(Anthropic's Claude and OpenAI's GPT) were used to assist with code development, refactoring, and
manuscript editing; all scientific content, results, and conclusions are my own and were checked
against the underlying data.
\end{acknowledgments}

\software{NumPy \citep{harris2020}, SciPy \citep{virtanen2020}, PyTorch \citep{paszke2019},
scikit-learn \citep{scikit-learn}, Matplotlib \citep{hunter2007},
VBBinaryLensing \citep{Bozza2018}, Astropy \citep{astropy2022}}

\facilities{AWS, Modal}

\section*{Data availability}
\sloppy     % long artifact paths; lets LaTeX stretch spaces rather than overflow the column
BinML, its trained weights and the full simulation-to-inference pipeline are openly available at
\url{https://github.com/kunalb541/BinML}. Every data-derived number, table and figure in this paper is generated by
the build script from the artifacts distributed with a release, in continuous integration from a clean checkout;
method and configuration constants, and the training settings and measurements of Appendix~\ref{app:infra}, are
typed. The frozen evaluation artifact (per-event logits, labels, weights, parameters and the test index, with a
hashed manifest) supports \S\ref{sec:results}--\ref{sec:limits}. The stress numbers re-reduce from their archived
per-event inputs alone (\code{validation/stress\_rescore\_archive/}); the floor, colour-calibration and every-class
stream results keep their per-event inputs in \code{validation/referee\_round\_archive/} but are reduced from the
runner's work directory (a reduction from the archive alone is not implemented); the full stress set's
$\bmlStressN$ million per-event predictions and the stage-5 weights are not distributed; and the no-blue-band slice
is read from the archived \code{metrics.json}. The RMDC26 release \citep{RGESPIT2026} is public (CC0) on the Hugging
Face Hub at \code{RGES-PIT/MachineLearning} (pinned revision \code{a338d5bab441b5caf551d2fea9469aadfdc81ec1});
\code{validation/gulls/rmdc26\_scores.csv.gz} holds every scored event's anomaly probability for every checkpoint,
from which the RMDC26 tables regenerate, and the recommended checkpoint is released as
\code{binml/weights/binml-gapaware.pt}. The scan, sub-day, relabelling, floor, between-season, occupancy,
noise-model, schedule and dataset-fact results rest on per-event data extracted from the release and regenerate by rerunning
their scripts on the pinned revision. Several archived runs have provenance limitations (uncommitted working trees
or missing hashes), recorded in the artifacts and in \code{validation/gulls/PROVENANCE.md}. The training set is
available on request and regenerates from the released simulation code; regenerated outside the original unpinned
cloud environment the data are close but not identical (for two held-out shards $\bmlRefMatchPct\%$ of the events
reproduce; the regenerated stress tiers label $\bmlStressNatAnomDeficit\%$ and $\bmlStressWidesepAnomDeficit\%$
fewer detectable anomalies), and the stress sweeps regenerate only with
\code{run\_shard --legacy-oor-mix --legacy-t0-pad}. Reproducing the checkpoints means rerunning that generation and
the training of Appendix~\ref{app:infra}.

\fussy
\appendix

\section{Simulation and training infrastructure}\label{app:infra}

The simulated data were generated on Amazon Web Services in the \texttt{\bmlSimRegion} region, mostly on
transient (Spot) instances of $\bmlSimVcpus$ virtual CPUs; the generator is embarrassingly parallel, one process
per self-contained shard seeded by its index, so a lost shard costs only itself. The same code, with the seed
bases of \S\ref{sec:limits}, produced the stress set. The fleets' software environment was not pinned, so a shard
regenerated elsewhere is close to, but not identical with, the original. We did not retain throughput or billing
records for these runs.

The released checkpoint is the last of $\bmlTrainStages$ stages, five of them warm-started from the previous one,
each keeping its epoch with the highest validation \texttt{NonPSPL} F1. Every stage used AdamW (weight decay
$0.01$), a one-cycle schedule (15\% warm-up, cosine decay) and cross-entropy with class-balanced weights (the
selection weights of \S\ref{sec:eval}, rescaled so that every class carries the same total weight). The base model
trained from scratch on $\bmlTrainBaseEvents$ natural-prior events (peak learning rate $4\times10^{-4}$, batch size
512, factor 2 on the anomaly class); stage 2 dropped that factor ($1.2\times10^{-4}$); stage 3 added the truncation
augmentation on half of all examples ($3\times10^{-5}$); stage 4 added enriched pools of 15 targeted regimes,
among them the low-$q$ and faint-source ($m = 23.5$--$25$) regimes ($5\times10^{-5}$, batch size 384 from here on);
stage 5 applied the detectability floor to truncated periodic classes; and stage 6 relabelled truncated binaries
as \texttt{PSPL} until their recorded onset and added pools of separations $s = 3$--$8$, timescales
$t_{\rm E} = 0.3$--$10$~d and further faint sources, $\bmlTrainFinalEvents$ million events in all, for
$\bmlTrainEpochs$ epochs; the released checkpoint is its zero-based epoch 4. Seeds are fixed ($\bmlSeedTrain$).
The fine-tunes of \S\ref{sec:limits} and \S\ref{sec:gulls} warm-start from the released weights or the
gap-augmented checkpoint with batch size 256 and the factor 2 on the anomaly class, at peak learning rates of
$5\times10^{-5}$ (colour calibration) and $10^{-4}$ (RMDC26). Every stage was trained on a personal computer
(\bmlTrainDevice) at $\bmlTrainEpsLo$--$\bmlTrainEpsHi$ training events per second including data loading,
augmentation and the validation pass; the $\bmlTrainStages$ stages took about $\bmlTrainHours$~hours of wall-clock,
about $\bmlTrainHoursTypical$ with a typical epoch in place of one stretched by a system sleep (which the throughput
range excludes). The compactness of the network is what keeps
training and survey-scale inference within reach of a single workstation.


\section{Training-seed spread and source-size dependence on RMDC26}\label{app:rmdc26}

Table~\ref{tab:rho} gives the single-lens false-alarm rate of every checkpoint of Table~\ref{tab:gulls} in six bins
of $\rho/|u_0|$, the ratio of source radius to impact parameter, which the finite-source fine-tunes of
\S\ref{sec:gulls:results} were designed to lower; Table~\ref{tab:seeds} gives the spread of three training seeds of
the recommended recipe, against which every difference between single runs in \S\ref{sec:gulls} is read.

\begin{deluxetable*}{lcccccc}
\tablecaption{Fraction of single lenses over the frozen threshold by $\rho/|u_0|$ bin, per simulated event, for the
checkpoints of Table~\ref{tab:gulls}. Generated from \code{validation/gulls/transfer\_tradeoff\_all.json}.\label{tab:rho}}
\tablehead{\colhead{checkpoint} & \colhead{$<0.03$} & \colhead{0.03--0.1} & \colhead{0.1--0.3} & \colhead{0.3--1} &
           \colhead{1--3} & \colhead{$> 3$}}
\startdata
shipped & 0.402 & 0.314 & 0.297 & 0.262 & 0.271 & 0.258 \\
gap augmentation (g08e12) & 0.046 & 0.100 & 0.229 & 0.429 & 0.646 & 0.674 \\
\quad + point-source control & 0.036 & 0.048 & 0.084 & 0.208 & 0.410 & 0.513 \\
\quad + finite source & 0.031 & 0.037 & 0.061 & 0.154 & 0.280 & 0.288 \\
\quad + finite source, pauses (rec.) & 0.033 & 0.050 & 0.097 & 0.224 & 0.330 & 0.333 \\
\qquad same recipe, seed 2 & 0.021 & 0.030 & 0.073 & 0.213 & 0.332 & 0.303 \\
\qquad same recipe, seed 3 & 0.034 & 0.058 & 0.132 & 0.276 & 0.414 & 0.375 \\

\enddata
\end{deluxetable*}

\begin{deluxetable*}{lcccccccccc}
\tabletypesize{\scriptsize}
\tablecaption{Differences between single training runs on RMDC26 against the spread of three seeds of the
recommended recipe: FA is the change in single-lens false-alarm rate at the frozen threshold (points), R$_x$ the
change in 1S2L recall at a single-lens budget of $x\%$, ``mean'' the change in mean recall over budgets up to
30\%; (w) columns weight by event rate; ``rec.'' is the recommended recipe. The first row is the range over its
three seeds; bold entries exceed it at the same budget and weighting (compared at full precision). The ``measured pauses'' row
compares the two schedule arms of Table~\ref{tab:gap} (another pool, from the shipped weights). Generated from \code{validation/gulls/transfer\_tradeoff\_all.json}.\label{tab:seeds}}
\tablehead{\colhead{comparison} & \colhead{FA} & \colhead{R$_{\bmlGullsBudgetLo}$} & \colhead{R$_{\bmlGullsBudgetMid}$} &
           \colhead{R$_{\bmlGullsBudgetHi}$} & \colhead{mean} & \colhead{FA (w)} & \colhead{R$_{\bmlGullsBudgetLo}$ (w)} &
           \colhead{R$_{\bmlGullsBudgetMid}$ (w)} & \colhead{R$_{\bmlGullsBudgetHi}$ (w)} & \colhead{mean (w)}}
\startdata
seed range (three seeds) & 2.7 & 0.041 & 0.042 & 0.029 & 0.017 & 3.9 & 0.016 & 0.022 & 0.028 & 0.015 \\
extra training (control $-$ gap aug.) & \textbf{-5.3} & \textbf{+0.059} & \textbf{+0.066} & +0.028 & +0.003 & \textbf{-4.9} & \textbf{+0.040} & \textbf{+0.038} & \textbf{+0.036} & +0.012 \\
finite source $-$ control & -1.5 & +0.021 & +0.030 & \textbf{+0.048} & \textbf{+0.049} & -1.0 & -0.015 & +0.010 & +0.012 & +0.015 \\
measured pauses $-$ random gaps & \textbf{-2.7} & \textbf{+0.058} & \textbf{+0.072} & \textbf{+0.058} & \textbf{+0.046} & -0.9 & \textbf{+0.038} & \textbf{+0.048} & \textbf{+0.034} & \textbf{+0.028} \\
recommended recipe, seed 1 $-$ finite source & +1.4 & +0.008 & +0.001 & -0.001 & -0.002 & +0.8 & \textbf{+0.024} & \textbf{+0.024} & +0.017 & \textbf{+0.016} \\
recommended recipe, seed 2 $-$ finite source & -0.2 & -0.033 & -0.041 & \textbf{-0.030} & \textbf{-0.019} & -1.7 & +0.008 & +0.002 & -0.011 & +0.000 \\
recommended recipe, seed 3 $-$ finite source & +2.5 & -0.032 & -0.032 & -0.019 & -0.010 & +2.2 & +0.012 & +0.002 & -0.004 & +0.007 \\

\enddata
\end{deluxetable*}

\bibliographystyle{aasjournalv7}
\bibliography{refs}

\end{document}
